% Requires in the preamble:
% \usepackage{amsmath,amssymb}
% \usepackage{tikz,pgfplots}
% \pgfplotsset{compat=1.18}

\chapter{Differential Equations: Laws Written as Derivatives}
\label{chap:differential-equations}

The last chapter was about approximation with responsibility.  We approximated
accumulated quantities, roots of equations, and derivatives from finite data.  Each
time, the question was not only what number an algorithm produced, but how much
that number could be trusted.

Differential equations carry the same lesson into a larger setting.

A differential equation does not usually give us the function directly.  It gives a law
for the function's derivative.  The unknown is not a number.  The unknown is a whole
function.

\[
\boxed{
\text{A differential equation is a law written in the language of change.}
}
\]

We have already seen hints of this idea.  Exponential growth was written as

\[
P'=kP.
\]

Cooling was written as

\[
T'=-k(T-T_s).
\]

Logistic growth was written as

\[
P'=rP\left(1-\frac{P}{K}\right).
\]

Those equations say how a quantity changes.  Solving them means finding the
function whose change obeys the law.

\section{Modeling with differential equations}
\label{sec:modeling-with-differential-equations}

A derivative measures an instant-by-instant rate.  A differential equation starts when a
situation gives that rate before it gives the function.

\[
\text{law of change}
\quad\Longrightarrow\quad
\text{equation for a derivative}
\quad\Longrightarrow\quad
\text{unknown function}.
\]

The direction is different from many earlier problems.  We are not given \(y(t)\) and
asked to find \(y'(t)\).  We are given information about \(y'(t)\) and asked to find
\(y(t)\).

\subsection{A law that gives a rate}
\label{subsec:law-gives-rate}

A hot cup of coffee cools quickly when it is much hotter than the room.  Later, when it
is only a little warmer than the room, it cools slowly.

That observation is the heart of Newton's law of cooling:

\[
\text{rate of cooling is proportional to the temperature difference}.
\]

Let

\[
T(t)=\text{temperature of the coffee after \(t\) minutes},
\]

and suppose the room temperature is

\[
T_s=20^\circ\text{C}.
\]

The temperature difference is

\[
T(t)-20.
\]

Cooling means this difference decreases.  A simple model is

\[
\boxed{
T'=-k(T-20),
}
\]

where \(k>0\) is a constant.

If \(k=0.08\), then

\[
T'=-0.08(T-20).
\]

When the coffee is \(80^\circ\text{C}\),

\[
T'=-0.08(80-20)=-4.8.
\]

So at that instant, the coffee is cooling at

\[
4.8^\circ\text{C per minute}.
\]

When the coffee is \(30^\circ\text{C}\),

\[
T'=-0.08(30-20)=-0.8.
\]

Now the cooling rate is only

\[
0.8^\circ\text{C per minute}.
\]

The equation captures the behavior:

\[
\text{large temperature difference} \quad\Longrightarrow\quad \text{fast cooling},
\]

\[
\text{small temperature difference} \quad\Longrightarrow\quad \text{slow cooling}.
\]

The constant \(k\) has units of \(1/\text{minute}\).  This is necessary because

\[
T'
\]

has units of degrees per minute, while

\[
T-20
\]

has units of degrees.

\begin{quote}
\textbf{A differential equation.}
A \emph{differential equation} is an equation involving an unknown function and one
or more of its derivatives.

The equation

\[
T'=-0.08(T-20)
\]

is a first-order differential equation because it involves the first derivative \(T'\), but
not higher derivatives such as \(T''\).
\end{quote}

The equation does not yet tell us the temperature at every time.  It tells us how the
temperature should be changing once the temperature is known.

That is the new kind of problem.

\subsection{Solution as an unknown function}
\label{subsec:solution-as-unknown-function}

A solution of a differential equation is a function that makes the equation true.

Start with a simpler equation:

\[
y'=1-y.
\]

This says that the rate of change of \(y\) depends on the current value of \(y\).

If \(y\) is below \(1\), then \(1-y>0\), so the solution should rise.

If \(y\) is above \(1\), then \(1-y<0\), so the solution should fall.

If \(y=1\), then \(1-y=0\), so the solution should stay flat.

Before solving the equation by a method, we can check proposed functions.

\paragraph{Example.}
Show that

\[
y(t)=1-e^{-t}
\]

is a solution of

\[
y'=1-y.
\]

Differentiate:

\[
y'(t)=e^{-t}.
\]

Now compute the right side of the differential equation:

\[
1-y(t)=1-\bigl(1-e^{-t}\bigr)=e^{-t}.
\]

The two sides agree:

\[
y'=e^{-t}
\qquad\text{and}\qquad
1-y=e^{-t}.
\]

Therefore

\[
\boxed{
y(t)=1-e^{-t}
}
\]

solves

\[
y'=1-y.
\]

\paragraph{Example.}
The function

\[
y(t)=1+2e^{-t}
\]

also solves the same equation.

Indeed,

\[
y'(t)=-2e^{-t},
\]

and

\[
1-y(t)=1-\bigl(1+2e^{-t}\bigr)=-2e^{-t}.
\]

So

\[
y'=1-y.
\]

The same differential equation has many solutions.  That should not be surprising.
An antiderivative also comes with a constant \(C\).  A differential equation often gives
a whole family of possible functions.

\begin{quote}
\textbf{Solution of a differential equation.}
A function \(y(t)\) is a \emph{solution} of a differential equation on an interval \(I\)
if \(y\) is differentiable on \(I\) and the equation is true for every \(t\) in \(I\).
\end{quote}

The interval matters.  A formula may solve an equation on one interval but fail to make
sense on another.  This is the same domain discipline we used for logarithms, square
roots, inverse trigonometric functions, and improper integrals.

\paragraph{A false proposed solution.}
Check whether

\[
y(t)=1+t
\]

solves

\[
y'=1-y.
\]

The derivative is

\[
y'(t)=1.
\]

But the right side is

\[
1-y(t)=1-(1+t)=-t.
\]

These are not equal for all \(t\).  They agree only when

\[
1=-t,
\]

or

\[
t=-1.
\]

A differential equation must hold on an interval, not just at one point.  Therefore

\[
\boxed{
y(t)=1+t \text{ is not a solution of } y'=1-y.
}
\]

\subsection{Initial conditions}
\label{subsec:initial-conditions}

A differential equation often describes a family of possible motions or histories.  A
starting value chooses one member of the family.

For

\[
y'=1-y,
\]

the functions

\[
y(t)=1+Ce^{-t}
\]

all solve the equation, for any constant \(C\).  Check:

\[
y'(t)=-Ce^{-t},
\]

and

\[
1-y(t)=1-\bigl(1+Ce^{-t}\bigr)=-Ce^{-t}.
\]

So

\[
y'=1-y.
\]

Now suppose we also know

\[
y(0)=0.
\]

Substitute \(t=0\):

\[
0=y(0)=1+Ce^0=1+C.
\]

Thus

\[
C=-1.
\]

So the solution with \(y(0)=0\) is

\[
\boxed{
y(t)=1-e^{-t}.
}
\]

If instead

\[
y(0)=2,
\]

then

\[
2=1+C,
\]

so

\[
C=1,
\]

and the solution is

\[
\boxed{
y(t)=1+e^{-t}.
}
\]

\begin{quote}
\textbf{Initial value problem.}
An \emph{initial value problem} combines a differential equation with a starting value.
For example,

\[
y'=1-y,
\qquad
y(0)=0.
\]

The differential equation gives the law of change.  The initial condition tells which
solution curve we are on.
\end{quote}

\begin{figure}[h]
\centering
\begin{tikzpicture}
\begin{axis}[
    width=0.82\textwidth,
    height=0.48\textwidth,
    xmin=0, xmax=4.2,
    ymin=-0.3, ymax=2.4,
    axis lines=left,
    xlabel={\(t\)},
    ylabel={\(y(t)\)},
    xtick={0,1,2,3,4},
    ytick={0,1,2},
    grid=both,
    grid style={line width=.1pt, draw=gray!25},
    legend style={draw=none, at={(0.98,0.98)}, anchor=north east}
]
\addplot[very thick, domain=0:4, samples=150] {1-exp(-x)};
\addlegendentry{\(y(0)=0\)}
\addplot[very thick, dashed, domain=0:4, samples=150] {1+exp(-x)};
\addlegendentry{\(y(0)=2\)}
\addplot[thick, dotted, domain=0:4, samples=2] {1};
\addlegendentry{\(y=1\)}
\addplot[only marks, mark=*] coordinates {(0,0) (0,2)};
\node[anchor=west] at (axis cs:2.15,1.07) {same equation};
\node[anchor=west] at (axis cs:2.15,0.82) {different initial values};
\end{axis}
\end{tikzpicture}
\caption{The equation \(y'=1-y\) gives a family of solution curves.  An initial condition selects one curve.}
\label{fig:initial-values-select-solutions}
\end{figure}

The picture shows what the algebra said.  Both solutions move toward \(y=1\).  One
starts below \(1\) and rises.  The other starts above \(1\) and falls.

\subsection{Checking a proposed solution}
\label{subsec:checking-proposed-solution}

Checking a solution is often easier than finding one.  It is also one of the best ways to
understand what the equation means.

\begin{quote}
\textbf{How to check a proposed solution.}

To check a proposed solution of a differential equation:

\begin{enumerate}
\item State the interval where the function is being considered.

\item Differentiate the proposed function.

\item Substitute the function and its derivative into the differential equation.

\item Check that the equation is true for every point of the interval.

\item If there is an initial condition, check it separately.
\end{enumerate}
\end{quote}

\paragraph{Example: checking a cooling model.}
Suppose the room temperature is \(20^\circ\text{C}\), and the proposed coffee
temperature is

\[
T(t)=20+60e^{-0.08t},
\qquad t\ge0.
\]

Check that it solves

\[
T'=-0.08(T-20),
\qquad
T(0)=80.
\]

First differentiate:

\[
T'(t)=60(-0.08)e^{-0.08t}=-4.8e^{-0.08t}.
\]

Now compute the right side:

\[
-0.08(T(t)-20)
=
-0.08\bigl(20+60e^{-0.08t}-20\bigr)
=
-4.8e^{-0.08t}.
\]

So

\[
T'=-0.08(T-20).
\]

Now check the initial condition:

\[
T(0)=20+60e^0=80.
\]

Therefore

\[
\boxed{
T(t)=20+60e^{-0.08t}
}
\]

solves the initial value problem.

\paragraph{What the formula says.}
The number \(20\) is the equilibrium temperature.  The term

\[
60e^{-0.08t}
\]

is the remaining temperature difference above the room.  It decays toward \(0\), so

\[
T(t)\to20
\]

as

\[
t\to\infty.
\]

The coffee approaches room temperature.  It does not pass below it in this model.

\begin{quote}
\textbf{A model check.}
A solution should satisfy the differential equation, the initial condition, and the
behavior of the situation.  Algebra alone is not enough if the model's interpretation
has been forgotten.
\end{quote}

Many differential equations cannot be solved by guessing a familiar formula.  But even
without a formula, the equation still gives information.  Since

\[
y'=F(t,y)
\]

tells the slope a solution must have at each point \((t,y)\), we can draw those slopes.
That picture is the next step.

\section{Direction fields}
\label{sec:direction-fields}

A first-order differential equation of the form

\[
y'=F(t,y)
\]

assigns a slope to each point in the \((t,y)\)-plane.  At the point \((t,y)\), a solution
curve passing through that point must have slope

\[
F(t,y).
\]

So even before solving the equation, we can draw a small line segment at many points.
The collection of these little segments is a picture of the differential equation.

\subsection{Slope field as a picture of a differential equation}
\label{subsec:slope-field-picture}

Return to the equation

\[
y'=1-y.
\]

At a point where \(y=0\),

\[
y'=1-0=1.
\]

So the solution curve has slope \(1\).

At a point where \(y=1\),

\[
y'=1-1=0.
\]

So the solution curve is horizontal.

At a point where \(y=2\),

\[
y'=1-2=-1.
\]

So the solution curve has slope \(-1\).

Notice that the right side does not depend on \(t\).  Every point along the same
horizontal line has the same slope.

\begin{figure}[h]
\centering
\begin{tikzpicture}
\begin{axis}[
    width=0.86\textwidth,
    height=0.52\textwidth,
    xmin=0, xmax=4.2,
    ymin=-0.2, ymax=2.4,
    axis lines=left,
    xlabel={\(t\)},
    ylabel={\(y\)},
    xtick={0,1,2,3,4},
    ytick={0,1,2},
    grid=both,
    grid style={line width=.1pt, draw=gray!25},
    legend style={draw=none, at={(0.98,0.97)}, anchor=north east}
]
% direction field for y' = 1-y
\foreach \xx in {0,0.5,...,4} {
  \foreach \yy in {0,0.25,...,2.25} {
    \pgfmathsetmacro{\m}{1-\yy}
    \pgfmathsetmacro{\len}{0.14}
    \pgfmathsetmacro{\dx}{\len/sqrt(1+\m*\m)}
    \pgfmathsetmacro{\dy}{\m*\dx}
    
    % Pre-calculate coordinate points to avoid parsing issues in axis cs
    \pgfmathsetmacro{\xOne}{\xx-\dx}
    \pgfmathsetmacro{\yOne}{\yy-\dy}
    \pgfmathsetmacro{\xTwo}{\xx+\dx}
    \pgfmathsetmacro{\yTwo}{\yy+\dy}
    
    % \edef expands the variables into hard numbers immediately
    \edef\temp{%
      \noexpand\draw[gray!65] (axis cs:\xOne,\yOne) -- (axis cs:\xTwo,\yTwo);%
    }
    % Execute the expanded command
    \temp
  }
}
% solution curves y=1+C e^{-t}
\addplot[very thick, domain=0:4, samples=150] {1-exp(-x)};
\addlegendentry{\(y(0)=0\)}
\addplot[very thick, dashed, domain=0:4, samples=150] {1+exp(-x)};
\addlegendentry{\(y(0)=2\)}
\addplot[thick, dotted, domain=0:4, samples=2] {1};
\addlegendentry{\(y=1\)}
\end{axis}
\end{tikzpicture}
\caption{A direction field draws the slope prescribed by \(y'=1-y\) at many points.  Solution curves follow the field.}
\label{fig:direction-field-one-minus-y}
\end{figure}

The little segments are not solution curves.  They are instructions for solution curves.
A true solution curve must pass through the field by matching the local direction at
each point.

\begin{quote}
\textbf{Direction field.}
For a differential equation

\[
y'=F(t,y),
\]

a \emph{direction field} or \emph{slope field} is a plot of small line segments, where
the segment at \((t,y)\) has slope \(F(t,y)\).
\end{quote}

The direction field is especially helpful when the equation is hard to solve exactly.  It
can show whether solutions rise or fall, whether they approach a stable value, and how
different initial conditions change the behavior.

\subsection{Equilibrium solutions}
\label{subsec:equilibrium-solutions}

Some solutions do not move at all.

For

\[
y'=1-y,
\]

the constant function

\[
y(t)=1
\]

is a solution.  Its derivative is

\[
y'=0,
\]

and the right side is

\[
1-y=1-1=0.
\]

So the equation is true.

This constant solution is called an equilibrium solution.

\begin{quote}
\textbf{Equilibrium solution.}
For an equation

\[
y'=F(t,y),
\]

a constant solution \(y(t)=c\) is called an \emph{equilibrium solution}.

If the equation is autonomous,

\[
y'=f(y),
\]

then equilibria occur at values \(c\) satisfying

\[
f(c)=0.
\]
\end{quote}

The word \emph{autonomous} means that the independent variable \(t\) does not appear
explicitly on the right side.  For example,

\[
y'=1-y
\]

is autonomous, while

\[
y'=t-y
\]

is not.

For an autonomous equation, all points with the same \(y\)-value have the same slope.
That is why the direction field in Figure \ref{fig:direction-field-one-minus-y} looks the
same across each horizontal row.

\subsection{Qualitative behavior}
\label{subsec:qualitative-behavior-direction-fields}

A direction field can tell us how solutions behave without giving exact formulas.

For

\[
y'=1-y,
\]

the sign of \(y'\) depends on whether \(y\) is below, equal to, or above \(1\).

\begin{center}
\begin{tabular}{c|c|c}
\textbf{region} & \textbf{sign of \(y'=1-y\)} & \textbf{behavior of solution} \\ \hline
\(y<1\) & positive & rises \\
\(y=1\) & zero & stays constant \\
\(y>1\) & negative & falls
\end{tabular}
\end{center}

So the line

\[
y=1
\]

acts like a attracting level.  Solutions below it rise toward it.  Solutions above it fall
toward it.

We can read this from the equation, from the sign table, or from the direction field.
The three views agree.

\paragraph{Example: logistic behavior without solving.}
Consider

\[
y'=y(1-y).
\]

This is the normalized logistic equation.  It appears when a quantity grows almost
exponentially when it is small, but slows as it approaches a carrying capacity.  Here the
carrying capacity has been scaled to \(1\).

Equilibrium solutions occur when

\[
y(1-y)=0.
\]

Thus

\[
y=0
\qquad\text{and}\qquad
y=1
\]

are equilibrium solutions.

Now read the sign of \(y(1-y)\).

\begin{center}
\begin{tabular}{c|c|c}
\textbf{region} & \textbf{sign of \(y(1-y)\)} & \textbf{behavior} \\ \hline
\(y<0\) & negative & falls \\
\(0<y<1\) & positive & rises \\
\(y>1\) & negative & falls
\end{tabular}
\end{center}

For a population model, we care mainly about

\[
y\ge0.
\]

If

\[
0<y(0)<1,
\]

the solution rises.

If

\[
y(0)>1,
\]

the solution falls.

In both cases, the solution is pushed toward

\[
y=1.
\]

The equilibrium \(y=1\) is stable.  The equilibrium \(y=0\) is unstable: a solution that
starts exactly at \(0\) stays there, but a solution that starts slightly above \(0\) moves
away from \(0\).

\begin{figure}[h]
\centering
\begin{tikzpicture}[xscale=1.2, yscale=1.2, every node/.style={font=\small}]
  \draw[->, thick] (0,-0.9) -- (0,2.0);
  \node[above] at (0,2.0) {\(y\)};

  \draw[thick] (-0.08,0) -- (0.08,0);
  \draw[thick] (-0.08,1) -- (0.08,1);

  \node[left] at (-0.08,0) {\(0\)};
  \node[left] at (-0.08,1) {\(1\)};

  % arrows
  \draw[->, thick] (0,-0.65) -- (0,-0.25);
  \draw[->, thick] (0,0.25) -- (0,0.65);
  \draw[->, thick] (0,1.65) -- (0,1.25);

  \node[right] at (0.18,-0.45) {\(y'<0\)};
  \node[right] at (0.18,0.45) {\(y'>0\)};
  \node[right] at (0.18,1.45) {\(y'<0\)};

  \node[right] at (0.18,0) {unstable equilibrium};
  \node[right] at (0.18,1) {stable equilibrium};
\end{tikzpicture}
\caption{For \(y'=y(1-y)\), the arrows show whether solutions move up or down.}
\label{fig:logistic-phase-line}
\end{figure}

This small vertical picture is a compressed direction field.  Because the equation is
autonomous, the sign depends only on \(y\), so one vertical line can summarize the
rise-and-fall behavior.

\subsection{Reading long-term behavior from the field}
\label{subsec:reading-long-term-behavior-field}

A direction field can answer questions that do not require an exact formula.

For

\[
y'=1-y,
\]

the field shows that solutions approach

\[
y=1
\]

as \(t\) increases.  The field also shows that the approach slows down, because slopes
near \(y=1\) are small.

For

\[
y'=y(1-y),
\]

the sign chart shows:

\[
0<y(0)<1
\quad\Longrightarrow\quad
\text{the solution rises toward }1,
\]

and

\[
y(0)>1
\quad\Longrightarrow\quad
\text{the solution falls toward }1.
\]

These conclusions are qualitative.  They tell the shape and long-term trend, not the
exact value at a specific time.

That is often enough for a first understanding of a model.  Before computing, we should
know what kind of behavior to expect.

\paragraph{Warning: the physical domain matters.}
The logistic equation

\[
y'=y(1-y)
\]

may be used for a normalized population.  In that interpretation,

\[
y<0
\]

does not make physical sense.  The equation still has mathematical behavior there, but
the model should not be interpreted as a population model for negative \(y\).

A differential equation is not the world.  It is a model of selected features of the world.

\paragraph{Warning: a direction field is not a uniqueness proof.}
A direction field can suggest that one solution curve passes through each point.  For
many equations in this chapter, that is true.  But it is not automatic.

Consider

\[
y'=\sqrt{|y|},
\qquad
y(0)=0.
\]

The constant function

\[
y(t)=0
\]

is a solution.

Another solution is

\[
y(t)=
\begin{cases}
0, & t\le 0,\\[3pt]
\dfrac{t^2}{4}, & t\ge 0.
\end{cases}
\]

For \(t>0\),

\[
y'(t)=\frac{t}{2}
\]

and

\[
\sqrt{|y(t)|}
=
\sqrt{\frac{t^2}{4}}
=
\frac{t}{2}.
\]

For \(t<0\), both sides are \(0\).  At \(t=0\), the derivative is also \(0\).  So this
piecewise function solves the same initial value problem.

Two different solution curves pass through the same initial point.

\begin{quote}
\textbf{What the warning teaches.}
A direction field is a powerful picture, but it is still a picture.  Clean uniqueness
requires hypotheses on the function \(F(t,y)\).  In this chapter, our examples will be
well behaved, but we should remember that the warning exists.
\end{quote}

The direction field gives the geometry of a differential equation.  The next question is
computational.  If the field tells us the slope at every point, can we follow those slopes
one small step at a time and build an approximate solution curve?

That step-by-step method is Euler's method.
\section{Euler's method}
\label{sec:eulers-method}

A direction field tells us the slope a solution must have at each point.  If we cannot
find a formula for the solution, we can still try to follow those slopes.

This is an old idea in new clothing.  The derivative gives a local linear approximation:

\[
y(x+h)\approx y(x)+y'(x)h.
\]

If the differential equation says

\[
y'=F(x,y),
\]

then at the point \((x,y)\), the slope is

\[
F(x,y).
\]

So one small step should look like

\[
y(x+h)\approx y(x)+hF(x,y).
\]

Euler's method is this local linear approximation repeated again and again.

\subsection{Following the slope field step by step}
\label{subsec:following-slope-field-step-by-step}

Start with the initial value problem

\[
y'=1-y,
\qquad
y(0)=0.
\]

We already know the direction field.  Below the line \(y=1\), slopes are positive, so
solutions rise.  Above \(y=1\), slopes are negative, so solutions fall.

Suppose we want an approximate value of \(y(2)\), but we pretend for the moment that
we do not know the exact formula for the solution.

Choose a step size

\[
h=0.5.
\]

At the starting point

\[
(x_0,y_0)=(0,0),
\]

the differential equation gives the slope

\[
y'=1-y=1-0=1.
\]

The local linear approximation says:

\[
y(0.5)\approx y(0)+0.5(1)=0.5.
\]

So our first approximate point is

\[
(x_1,y_1)=(0.5,0.5).
\]

Now repeat the same idea.  At \((0.5,0.5)\), the slope is

\[
1-y_1=1-0.5=0.5.
\]

Therefore

\[
y_2=y_1+h(1-y_1)=0.5+0.5(0.5)=0.75.
\]

So

\[
(x_2,y_2)=(1,0.75).
\]

Continue:

\[
y_3=0.75+0.5(1-0.75)=0.875,
\]

and

\[
y_4=0.875+0.5(1-0.875)=0.9375.
\]

Thus Euler's method gives

\[
\boxed{
y(2)\approx0.9375.
}
\]

The computation is best seen in a table.

\begin{center}
\begin{tabular}{c|c|c|c|c}
\textbf{step} & \(x_k\) & \(y_k\) & slope \(1-y_k\) & next value \(y_{k+1}\) \\ \hline
\(0\) & \(0\) & \(0\) & \(1\) & \(0+0.5(1)=0.5\) \\
\(1\) & \(0.5\) & \(0.5\) & \(0.5\) & \(0.5+0.5(0.5)=0.75\) \\
\(2\) & \(1.0\) & \(0.75\) & \(0.25\) & \(0.75+0.5(0.25)=0.875\) \\
\(3\) & \(1.5\) & \(0.875\) & \(0.125\) & \(0.875+0.5(0.125)=0.9375\)
\end{tabular}
\end{center}

Each row says:

\[
\text{current height}+\text{step size}\cdot\text{current slope}
=
\text{next height}.
\]

\begin{figure}[h]
\centering
\begin{tikzpicture}
\begin{axis}[
    width=0.86\textwidth,
    height=0.52\textwidth,
    xmin=0, xmax=2.2,
    ymin=-0.1, ymax=1.2,
    axis lines=left,
    xlabel={\(x\)},
    ylabel={\(y\)},
    xtick={0,0.5,1,1.5,2},
    ytick={0,0.5,1},
    grid=both,
    grid style={line width=.1pt, draw=gray!25},
    legend style={draw=none, at={(0.98,0.25)}, anchor=south east}
]
% direction field for y' = 1-y
% direction field for y' = 1-y
\foreach \xx in {0,0.5,...,4} {
  \foreach \yy in {0,0.25,...,2.25} {
    \pgfmathsetmacro{\m}{1-\yy}
    \pgfmathsetmacro{\len}{0.14}
    \pgfmathsetmacro{\dx}{\len/sqrt(1+\m*\m)}
    \pgfmathsetmacro{\dy}{\m*\dx}
    
    % Pre-calculate coordinates
    \pgfmathsetmacro{\xOne}{\xx-\dx}
    \pgfmathsetmacro{\yOne}{\yy-\dy}
    \pgfmathsetmacro{\xTwo}{\xx+\dx}
    \pgfmathsetmacro{\yTwo}{\yy+\dy}
    
    % \edef forces expansion of the variables into raw numbers immediately
    \edef\temp{%
      \noexpand\draw[gray!65] (axis cs:\xOne,\yOne) -- (axis cs:\xTwo,\yTwo);%
    }
    % Execute the expanded command, passing hard coordinates to the axis
    \temp
  }
}

% exact solution
\addplot[very thick, domain=0:2.05, samples=150] {1-exp(-x)};
\addlegendentry{exact \(y=1-e^{-x}\)}

% Euler broken line
\addplot[very thick, dashed, mark=*] coordinates {
  (0,0) (0.5,0.5) (1,0.75) (1.5,0.875) (2,0.9375)
};
\addlegendentry{Euler, \(h=0.5\)}

% equilibrium
\addplot[thick, dotted, domain=0:2.05, samples=2] {1};
\node[anchor=west] at (axis cs:0.08,1.04) {\(y=1\)};
\end{axis}
\end{tikzpicture}
\caption{Euler's method follows the slope field by taking one tangent-line step at a time.}
\label{fig:euler-y-prime-one-minus-y}
\end{figure}

In this example, we can compare with the exact solution.  The initial value problem

\[
y'=1-y,
\qquad
y(0)=0
\]

has solution

\[
y=1-e^{-x}.
\]

Therefore

\[
y(2)=1-e^{-2}\approx0.8647.
\]

Euler's method with step size \(h=0.5\) gave

\[
0.9375.
\]

So it overestimates the true value.  The reason is visible in the picture: the solution is
concave down while it rises toward \(1\), so tangent-line steps tend to sit above the
curve.

\begin{quote}
\textbf{Euler's method.}
For an initial value problem

\[
y'=F(x,y),
\qquad
y(x_0)=y_0,
\]

choose a step size \(h\).  Define

\[
x_{k+1}=x_k+h
\]

and

\[
\boxed{
y_{k+1}=y_k+hF(x_k,y_k).
}
\]

The points

\[
(x_0,y_0),\ (x_1,y_1),\ (x_2,y_2),\ldots
\]

approximate points on the solution curve.
\end{quote}

Euler's method is simple because it uses only the slope at the beginning of each step.
That simplicity is also its weakness.

\subsection{Step size and accumulated error}
\label{subsec:step-size-accumulated-error}

Euler's method makes a small error at each step.  Then it uses the approximate point
from that step to compute the next slope.  So errors do not merely appear; they can
accumulate.

For the same problem,

\[
y'=1-y,
\qquad
y(0)=0,
\]

we can approximate \(y(2)\) with several step sizes.

\[
\text{exact value}=1-e^{-2}\approx0.864665.
\]

\begin{center}
\begin{tabular}{c|c|c|c}
\textbf{step size} & \textbf{number of steps} & \textbf{Euler approximation to \(y(2)\)} & \textbf{absolute error} \\ \hline
\(1\) & \(2\) & \(1.000000\) & \(0.135335\) \\
\(0.5\) & \(4\) & \(0.937500\) & \(0.072835\) \\
\(0.25\) & \(8\) & \(0.899887\) & \(0.035222\) \\
\(0.1\) & \(20\) & \(0.878423\) & \(0.013759\) \\
\(0.05\) & \(40\) & \(0.871488\) & \(0.006823\)
\end{tabular}
\end{center}

The smaller step sizes give better approximations here.  This is what we hope for:
more, shorter tangent-line steps should trace the solution curve more faithfully.

But there is a cost.  Smaller steps require more computation.

\[
h=0.5
\quad\Longrightarrow\quad
4\text{ steps to reach }x=2,
\]

while

\[
h=0.05
\quad\Longrightarrow\quad
40\text{ steps to reach }x=2.
\]

The tradeoff is familiar from numerical integration:

\[
\text{smaller step size}
\quad\Longrightarrow\quad
\text{usually smaller error, but more work}.
\]

The word ``usually'' belongs in the sentence.  A numerical method is not magic.  If the
step size is too large, the method may follow the wrong behavior.

\paragraph{A large-step warning.}
Consider

\[
y'=-4y,
\qquad
y(0)=1.
\]

The exact solution is

\[
y=e^{-4x}.
\]

It stays positive and decays toward \(0\).

Euler's method gives

\[
y_{k+1}=y_k+h(-4y_k)=(1-4h)y_k.
\]

If

\[
h=1,
\]

then

\[
y_{k+1}=-3y_k.
\]

The Euler values are

\[
1,\ -3,\ 9,\ -27,\ldots
\]

They alternate signs and grow in size.  That is the opposite of the true solution.

The differential equation describes stable decay.  The numerical method, with a large
step, creates fake instability.

\begin{quote}
\textbf{Numerical warning.}
A differential equation and a numerical method are different objects.  The equation may
have one behavior, while a poor step size in the method can produce another.
\end{quote}

This does not make Euler's method useless.  It makes it honest.  We should use it with
step-size discipline and with checks against qualitative behavior.

\subsection{Improved Euler method}
\label{subsec:improved-euler-method}

Euler's method uses the slope at the beginning of the step.  But if the slope changes
during the step, the beginning slope may be a poor representative of the whole interval.

A natural improvement is to use two slopes:

\[
\text{beginning slope}
\qquad\text{and}\qquad
\text{estimated ending slope}.
\]

Then average them.

Start at

\[
(x_k,y_k).
\]

Let

\[
k_1=F(x_k,y_k)
\]

be the beginning slope.  Euler's ordinary prediction is

\[
\widetilde y_{k+1}=y_k+hk_1.
\]

This gives a predicted endpoint

\[
(x_k+h,\widetilde y_{k+1}).
\]

Now compute the slope there:

\[
k_2=F(x_k+h,\widetilde y_{k+1}).
\]

The improved Euler step uses the average slope:

\[
\boxed{
y_{k+1}
=
y_k+\frac{h}{2}(k_1+k_2).
}
\]

\begin{quote}
\textbf{Improved Euler method.}
For

\[
y'=F(x,y),
\qquad
y(x_0)=y_0,
\]

set

\[
k_1=F(x_k,y_k),
\]

\[
\widetilde y_{k+1}=y_k+hk_1,
\]

\[
k_2=F(x_k+h,\widetilde y_{k+1}),
\]

and

\[
\boxed{
y_{k+1}=y_k+\frac{h}{2}(k_1+k_2).
}
\]
\end{quote}

This method is also called Heun's method.  It is not the most powerful numerical method
for differential equations, but it fixes one clear defect of Euler's method: it does not rely
only on the starting slope.

\paragraph{Example.}
Use improved Euler with \(h=0.5\) for

\[
y'=1-y,
\qquad
y(0)=0.
\]

The computations are:

\begin{center}
\begin{tabular}{c|c|c|c|c|c}
\textbf{step} & \(x_k\) & \(y_k\) & \(k_1\) & \(k_2\) & \(y_{k+1}\) \\ \hline
\(0\) & \(0\) & \(0\) & \(1\) & \(0.5\) & \(0.375\) \\
\(1\) & \(0.5\) & \(0.375\) & \(0.625\) & \(0.3125\) & \(0.609375\) \\
\(2\) & \(1.0\) & \(0.609375\) & \(0.390625\) & \(0.1953125\) & \(0.755859\) \\
\(3\) & \(1.5\) & \(0.755859\) & \(0.244141\) & \(0.122070\) & \(0.847412\)
\end{tabular}
\end{center}

So

\[
y(2)\approx0.847412.
\]

The exact value is

\[
0.864665.
\]

Euler's method with the same step size gave

\[
0.937500.
\]

In this example, improved Euler is much closer.

The price is also clear.  Ordinary Euler evaluates \(F\) once per step.  Improved Euler
evaluates \(F\) twice per step.  Better information usually costs more computation.

\subsection{Numerical solutions as approximate functions}
\label{subsec:numerical-solutions-approximate-functions}

Euler's method does not produce a formula like

\[
y=1-e^{-x}.
\]

It produces a list of points:

\[
(x_0,y_0),\ (x_1,y_1),\ (x_2,y_2),\ldots.
\]

If we connect those points by straight line segments, we get a piecewise linear
approximation to the solution curve.  So Euler's method approximates not only one
number, but a whole function.

This distinction matters.

A numerical answer such as

\[
y(2)\approx0.9375
\]

is not complete unless we know how it was produced.  We should report:

\[
\text{the method},
\qquad
\text{the step size},
\qquad
\text{the interval},
\qquad
\text{and any available error information}.
\]

For example,

\[
\boxed{
\text{Euler's method with \(h=0.5\) gives \(y(2)\approx0.9375\).}
}
\]

That statement is honest.  It names the approximation.

\begin{quote}
\textbf{What Euler's method teaches.}

A differential equation gives local instructions.  Euler's method turns those instructions
into a path by repeating one simple idea:

\[
\text{new value}=\text{old value}+\text{slope}\cdot\text{step}.
\]
\end{quote}

This is enough when no exact formula is available.  But sometimes the equation itself
contains a hidden integral.  In those cases, we can do better than following short
tangent steps.  We can solve the equation exactly.

The first such equations are separable.

\section{Separable equations}
\label{sec:separable-equations}

Euler's method follows a differential equation numerically.  A separable equation can
sometimes be solved by integration.

The idea is simple.  Some differential equations can be rearranged so that all the \(y\)'s
are on one side and all the \(x\)'s are on the other.

Then the equation becomes an integral statement.

\[
\text{separate}
\quad\Longrightarrow\quad
\text{integrate}
\quad\Longrightarrow\quad
\text{solve for the constant}.
\]

\subsection{Separating variables}
\label{subsec:separating-variables}

Start with the differential equation

\[
\frac{dy}{dx}=xy.
\]

This says that the rate of change of \(y\) is the product of the current input \(x\) and
the current output \(y\).

If \(y\neq0\), we can divide by \(y\):

\[
\frac{1}{y}\frac{dy}{dx}=x.
\]

In differential notation, this is written

\[
\frac{dy}{y}=x\,dx.
\]

Now integrate both sides:

\[
\int \frac{1}{y}\,dy=\int x\,dx.
\]

Thus

\[
\ln|y|=\frac{x^2}{2}+C.
\]

Exponentiating gives

\[
|y|=e^C e^{x^2/2}.
\]

The sign of \(y\) can be absorbed into the constant, so the nonzero solutions are

\[
y=Ce^{x^2/2}.
\]

Now apply an initial condition.  If

\[
y(0)=2,
\]

then

\[
2=Ce^0,
\]

so

\[
C=2.
\]

Therefore

\[
\boxed{
y=2e^{x^2/2}.
}
\]

Check the answer:

\[
y'=2e^{x^2/2}\cdot x
=
xy.
\]

The formula satisfies the differential equation and the initial condition.

\begin{quote}
\textbf{Separable equation.}
A first-order differential equation is \emph{separable} if it can be written in the form

\[
\frac{dy}{dx}=g(x)h(y).
\]

On intervals where \(h(y)\neq0\), we can write

\[
\frac{1}{h(y)}\,dy=g(x)\,dx
\]

and integrate:

\[
\boxed{
\int \frac{1}{h(y)}\,dy=\int g(x)\,dx+C.
}
\]
\end{quote}

The separation notation is compact, but it is not meaningless symbol pushing.  It is the
chain rule in disguise.

Suppose

\[
H'(y)=\frac{1}{h(y)}
\]

and

\[
G'(x)=g(x).
\]

If \(y=y(x)\) satisfies

\[
y'=g(x)h(y),
\]

then

\[
\frac{d}{dx}H(y(x))
=
H'(y)y'
=
\frac{1}{h(y)}g(x)h(y)
=
g(x)
=
G'(x).
\]

Therefore

\[
H(y(x))=G(x)+C.
\]

That is exactly what separation and integration produced.

\paragraph{Workflow.}
To solve a separable equation:

\begin{enumerate}
\item Write it in the form

\[
\frac{dy}{dx}=g(x)h(y).
\]

\item Move all \(y\)-factors to the left and all \(x\)-factors to the right:

\[
\frac{1}{h(y)}\,dy=g(x)\,dx.
\]

\item Integrate both sides.

\item Use the initial condition to determine the constant.

\item Solve explicitly for \(y\), if possible.

\item Check for constant solutions lost by division.
\end{enumerate}

The last step is easy to forget and important.

\paragraph{Warning: division can lose solutions.}
Consider

\[
y'=y(1-y).
\]

This equation is separable.  But if we divide by

\[
y(1-y),
\]

we have assumed

\[
y\neq0
\qquad\text{and}\qquad
y\neq1.
\]

The constant functions

\[
y=0
\qquad\text{and}\qquad
y=1
\]

are both solutions, because in each case

\[
y'=0
\]

and

\[
y(1-y)=0.
\]

If we divide by \(y(1-y)\) before noticing them, we may lose them.

\begin{quote}
\textbf{Separation rule of honesty.}
When you divide by an expression involving \(y\), separately check the values of \(y\)
that make that expression zero.  They often give equilibrium solutions.
\end{quote}

This warning will matter even more for logistic models in the next section.

\subsection{Growth and decay revisited}
\label{subsec:growth-decay-revisited-separable}

The most famous separable equation is

\[
y'=ky,
\]

where \(k\) is constant.

This says:

\[
\text{rate of change is proportional to current amount}.
\]

Separate variables:

\[
\frac{dy}{y}=k\,dt,
\qquad y\neq0.
\]

Integrate:

\[
\ln|y|=kt+C.
\]

Exponentiate:

\[
y=Ce^{kt}.
\]

If

\[
y(0)=y_0,
\]

then

\[
C=y_0.
\]

So

\[
\boxed{
y(t)=y_0e^{kt}.
}
\]

The formula also works when \(y_0=0\), giving the equilibrium solution

\[
y(t)=0.
\]

The sign of \(k\) determines the long-term behavior.

\begin{center}
\begin{tabular}{c|c|c}
\textbf{equation} & \textbf{condition} & \textbf{behavior} \\ \hline
\(y'=ky\) & \(k>0\) & exponential growth \\
\(y'=ky\) & \(k<0\) & exponential decay \\
\(y'=ky\) & \(k=0\) & constant solution
\end{tabular}
\end{center}

\paragraph{Example: radioactive decay.}
A radioactive sample has mass \(M(t)\), measured in grams, and satisfies

\[
M'=-0.03M,
\qquad
M(0)=80.
\]

The solution is

\[
M(t)=80e^{-0.03t}.
\]

The mass decays toward \(0\), but never reaches \(0\) in finite time.

To find the half-life, solve

\[
40=80e^{-0.03t}.
\]

Divide by \(80\):

\[
\frac12=e^{-0.03t}.
\]

Take logarithms:

\[
\ln\left(\frac12\right)=-0.03t.
\]

Since

\[
\ln\left(\frac12\right)=-\ln2,
\]

we get

\[
t=\frac{\ln2}{0.03}.
\]

Thus

\[
\boxed{
t\approx23.1.
}
\]

The half-life is about \(23.1\) time units.

\paragraph{Example: variable relative growth.}
Solve

\[
y'=ty,
\qquad
y(0)=3.
\]

Separate:

\[
\frac{dy}{y}=t\,dt.
\]

Integrate:

\[
\ln|y|=\frac{t^2}{2}+C.
\]

Since \(y(0)=3\), we get

\[
C=\ln3.
\]

Thus

\[
\boxed{
y=3e^{t^2/2}.
}
\]

Here the relative growth rate is not constant.  It is \(t\).  The exponent is therefore the
accumulated growth rate:

\[
\int_0^t s\,ds=\frac{t^2}{2}.
\]

This is a useful way to remember the answer:

\[
\boxed{
\text{amount}=\text{initial amount}\cdot e^{\text{accumulated relative rate}}.
}
\]

\subsection{Cooling and mixing}
\label{subsec:cooling-and-mixing}

Many physical models become separable because the rate depends on a difference from
equilibrium.

Newton's law of cooling says

\[
T'=-k(T-T_s),
\qquad k>0,
\]

where \(T_s\) is the surrounding temperature.

Separate variables:

\[
\frac{dT}{T-T_s}=-k\,dt.
\]

Integrate:

\[
\ln|T-T_s|=-kt+C.
\]

Exponentiate:

\[
T-T_s=Ce^{-kt}.
\]

If

\[
T(0)=T_0,
\]

then

\[
C=T_0-T_s.
\]

So

\[
\boxed{
T(t)=T_s+(T_0-T_s)e^{-kt}.
}
\]

This single formula describes cooling and warming.

If \(T_0>T_s\), then \(T_0-T_s>0\), and the object cools down toward \(T_s\).

If \(T_0<T_s\), then \(T_0-T_s<0\), and the object warms up toward \(T_s\).

The exponential term decays to \(0\), so in both cases

\[
T(t)\to T_s
\qquad\text{as}\qquad
t\to\infty.
\]

\paragraph{Example: coffee cooling.}
A cup of coffee is initially

\[
80^\circ\text{C}.
\]

The room is

\[
20^\circ\text{C}.
\]

Suppose

\[
T'=-0.08(T-20),
\qquad
T(0)=80.
\]

Then

\[
T(t)=20+(80-20)e^{-0.08t}.
\]

Thus

\[
\boxed{
T(t)=20+60e^{-0.08t}.
}
\]

When will the coffee reach \(40^\circ\text{C}\)?

Solve

\[
40=20+60e^{-0.08t}.
\]

Then

\[
20=60e^{-0.08t},
\]

so

\[
\frac13=e^{-0.08t}.
\]

Take logarithms:

\[
-\ln3=-0.08t.
\]

Therefore

\[
t=\frac{\ln3}{0.08}\approx13.7.
\]

The coffee reaches \(40^\circ\text{C}\) after about

\[
\boxed{13.7\text{ minutes}.}
\]

\begin{figure}[h]
\centering
\begin{tikzpicture}
\begin{axis}[
    width=0.82\textwidth,
    height=0.48\textwidth,
    xmin=0, xmax=35,
    ymin=15, ymax=85,
    axis lines=left,
    xlabel={\(t\) (minutes)},
    ylabel={\(T(t)\) (\(^\circ\)C)},
    xtick={0,10,20,30},
    ytick={20,40,60,80},
    grid=both,
    grid style={line width=.1pt, draw=gray!25},
    legend style={draw=none, at={(0.98,0.98)}, anchor=north east}
]
\addplot[very thick, domain=0:35, samples=200] {20+60*exp(-0.08*x)};
\addlegendentry{\(T=20+60e^{-0.08t}\)}
\addplot[thick, dotted, domain=0:35, samples=2] {20};
\addlegendentry{\(T=20\)}
\addplot[thick, dashed, domain=0:35, samples=2] {40};
\addplot[only marks, mark=*] coordinates {(13.7327,40)};
\node[anchor=south west] at (axis cs:13.7327,40) {\(t\approx13.7\)};
\end{axis}
\end{tikzpicture}
\caption{Newton's law of cooling gives exponential approach to the surrounding temperature.}
\label{fig:cooling-separable}
\end{figure}

Mixing can produce the same kind of equation.

\paragraph{Example: flushing salt from a tank.}
A tank contains \(100\) liters of well-mixed salt water.  Initially there are \(10\) grams of
salt in the tank.  Pure water flows in at \(5\) liters per minute, and the well-mixed
solution flows out at the same rate.

Let

\[
S(t)=\text{grams of salt in the tank after \(t\) minutes}.
\]

Since the volume stays \(100\) liters, the concentration in the tank is

\[
\frac{S(t)}{100}
\]

grams per liter.

Pure water brings in no salt.  Salt leaves at the rate

\[
5\cdot\frac{S(t)}{100}
=
\frac{S(t)}{20}
\]

grams per minute.

Therefore

\[
S'=-\frac{S}{20}.
\]

This is separable:

\[
\frac{dS}{S}=-\frac{1}{20}\,dt.
\]

Integrate:

\[
\ln|S|=-\frac{t}{20}+C.
\]

Since \(S(0)=10\), we get

\[
\boxed{
S(t)=10e^{-t/20}.
}
\]

The salt amount decays exponentially.

When will only \(1\) gram remain?

\[
1=10e^{-t/20}.
\]

So

\[
\frac{1}{10}=e^{-t/20}.
\]

Taking logarithms gives

\[
-\ln10=-\frac{t}{20}.
\]

Thus

\[
\boxed{
t=20\ln10\approx46.1\text{ minutes}.
}
\]

The tank never becomes perfectly salt-free in finite time, but the amount of salt can be
made as small as desired.

\begin{quote}
\textbf{Equilibrium difference pattern.}
Equations of the form

\[
Q'=-k(Q-Q_{\text{eq}})
\]

have solutions

\[
\boxed{
Q(t)=Q_{\text{eq}}+\bigl(Q(0)-Q_{\text{eq}}\bigr)e^{-kt}.
}
\]

The quantity \(Q\) approaches the equilibrium value \(Q_{\text{eq}}\) exponentially.
\end{quote}

Cooling, charging circuits, simple mixing, and many adjustment models share this same
mathematical skeleton.

\subsection{Implicit solutions}
\label{subsec:implicit-solutions-separable}

Separation does not always lead to a clean formula

\[
y=\text{something involving }x.
\]

Sometimes the best answer is an equation that \(x\) and \(y\) satisfy together.  Such an
answer is called an implicit solution.

Consider

\[
y'=\frac{x}{1+y^2},
\qquad
y(0)=0.
\]

Separate variables:

\[
(1+y^2)\,dy=x\,dx.
\]

Integrate:

\[
\int(1+y^2)\,dy=\int x\,dx.
\]

So

\[
y+\frac{y^3}{3}=\frac{x^2}{2}+C.
\]

Use the initial condition \(y(0)=0\):

\[
0+\frac{0^3}{3}=0+C,
\]

so

\[
C=0.
\]

Thus the solution is given implicitly by

\[
\boxed{
y+\frac{y^3}{3}=\frac{x^2}{2}.
}
\]

This equation is a perfectly meaningful solution.  It tells us which points \((x,y)\)
lie on the solution curve.

We can also check it without solving for \(y\).  Differentiate both sides with respect to
\(x\):

\[
\frac{d}{dx}\left(y+\frac{y^3}{3}\right)
=
\frac{d}{dx}\left(\frac{x^2}{2}\right).
\]

By the chain rule,

\[
y'+y^2y'=x.
\]

Factor:

\[
(1+y^2)y'=x.
\]

Therefore

\[
y'=\frac{x}{1+y^2}.
\]

So the implicit equation really does describe a solution of the differential equation.

\begin{quote}
\textbf{Implicit solution.}
An \emph{implicit solution} of a differential equation is an equation involving \(x\)
and \(y\) whose differentiable branches satisfy the differential equation.

An explicit solution gives \(y\) directly as a function of \(x\).  An implicit solution may
leave \(x\) and \(y\) tied together.
\end{quote}

\paragraph{Example: a family of hyperbolas.}
Solve

\[
y'=\frac{x}{y},
\qquad
y(0)=2.
\]

Separate:

\[
y\,dy=x\,dx.
\]

Integrate:

\[
\frac{y^2}{2}=\frac{x^2}{2}+C.
\]

Multiply by \(2\):

\[
y^2=x^2+C_1.
\]

Use the initial condition:

\[
2^2=0^2+C_1,
\]

so

\[
C_1=4.
\]

Thus

\[
\boxed{
y^2-x^2=4.
}
\]

This implicit equation has two branches:

\[
y=\sqrt{x^2+4}
\]

and

\[
y=-\sqrt{x^2+4}.
\]

The initial condition \(y(0)=2\) selects the upper branch:

\[
\boxed{
y=\sqrt{x^2+4}.
}
\]

The differential equation

\[
y'=\frac{x}{y}
\]

does not make sense when \(y=0\).  The selected branch never reaches \(y=0\), so the
solution stays safely inside the domain of the equation.

\paragraph{Domains still matter.}
Solving a differential equation is not just manipulating symbols until a formula appears.
A solution must live on an interval where the function is differentiable and the equation
makes sense.

For

\[
y'=\frac{x}{y},
\]

any solution must have

\[
y\neq0.
\]

For

\[
y'=\frac{1}{x}y,
\]

any solution interval must avoid

\[
x=0.
\]

A formula without a domain is not yet a complete solution.

\begin{quote}
\textbf{What separable equations give us.}
Separation turns some differential equations into integral problems.  It connects this
chapter back to the Fundamental Theorem:

\[
\text{a law for a derivative}
\quad\Longrightarrow\quad
\text{an integral relation}
\quad\Longrightarrow\quad
\text{a solution curve}.
\]
\end{quote}

The next separable model has two equilibrium solutions instead of one.  It grows almost
exponentially at first, slows as it approaches a limit, and has a characteristic
slow--fast--slow shape.

That model is the logistic equation.
\section{Logistic and threshold models}
\label{sec:logistic-threshold-models}

Separable equations gave us exact formulas when the variables could be pulled apart.
But the formula is not always the first thing we should look for.

For an autonomous equation

\[
y'=f(y),
\]

the graph of \(f\) already tells a story.  It tells where solutions rise, where they fall,
and which equilibrium values attract nearby solutions.

The logistic equation is the best first example.  It can be solved by separation, but it is
often understood first from its phase line.

\subsection{Carrying capacity}
\label{subsec:carrying-capacity}

Pure exponential growth says

\[
P'=rP,
\qquad r>0.
\]

The rate is proportional to the current population.  That can describe early growth, when
resources are plentiful.  It cannot describe growth forever.

A population in a fixed environment has limited food, space, or available hosts.  The
model should slow down as the population approaches a maximum sustainable size.

Let

\[
K=\text{carrying capacity}.
\]

The simplest slowing factor is

\[
1-\frac{P}{K}.
\]

When

\[
P\ll K,
\]

this factor is close to \(1\), so the equation looks almost exponential:

\[
P'\approx rP.
\]

When

\[
P=K,
\]

the factor becomes \(0\), so growth stops.

This gives the logistic equation:

\[
\boxed{
P'=rP\left(1-\frac{P}{K}\right),
\qquad r>0,\quad K>0.
}
\]

The two factors have different meanings:

\[
rP
\]

says that a larger current population creates more births, while

\[
1-\frac{P}{K}
\]

says that limited remaining capacity slows the growth.

\paragraph{Solving the logistic equation.}
Assume for the moment that

\[
P\neq0
\qquad\text{and}\qquad
P\neq K.
\]

Then we can separate variables:

\[
\frac{dP}{P\left(1-\frac{P}{K}\right)}=r\,dt.
\]

Since

\[
P\left(1-\frac{P}{K}\right)=\frac{P(K-P)}{K},
\]

we have

\[
\frac{1}{P\left(1-\frac{P}{K}\right)}
=
\frac{K}{P(K-P)}.
\]

Use partial fractions:

\[
\frac{K}{P(K-P)}
=
\frac1P+\frac1{K-P}.
\]

Therefore

\[
\int\left(\frac1P+\frac1{K-P}\right)\,dP
=
\int r\,dt.
\]

The left side is

\[
\ln|P|-\ln|K-P|.
\]

So

\[
\ln\left|\frac{P}{K-P}\right|=rt+C.
\]

Exponentiating gives

\[
\frac{P}{K-P}=Ae^{rt},
\]

where \(A\) is a nonzero constant.  Solving for \(P\),

\[
P=Ae^{rt}(K-P),
\]

so

\[
P(1+Ae^{rt})=KAe^{rt}.
\]

Thus

\[
P(t)=\frac{KAe^{rt}}{1+Ae^{rt}}.
\]

Dividing numerator and denominator by \(Ae^{rt}\), we get the usual form:

\[
\boxed{
P(t)=\frac{K}{1+Be^{-rt}}.
}
\]

If

\[
P(0)=P_0,
\qquad P_0>0,
\]

then

\[
P_0=\frac{K}{1+B}.
\]

So

\[
1+B=\frac{K}{P_0},
\]

and therefore

\[
\boxed{
B=\frac{K-P_0}{P_0}.
}
\]

Thus, for \(P_0>0\),

\[
\boxed{
P(t)=\frac{K}{1+\left(\frac{K-P_0}{P_0}\right)e^{-rt}}.
}
\]

This formula includes the case \(P_0=K\), where the constant in the denominator is
\(0\), and the solution is

\[
P(t)=K.
\]

The solution \(P(t)=0\) is also an equilibrium solution, but it was lost when we divided
by \(P\).  This is the same warning we saw in separable equations: division can erase
constant solutions.

\paragraph{Example.}
A population grows according to

\[
P'=0.4P\left(1-\frac{P}{1000}\right),
\qquad
P(0)=100.
\]

Here

\[
r=0.4,
\qquad
K=1000,
\qquad
P_0=100.
\]

So

\[
B=\frac{1000-100}{100}=9.
\]

The solution is

\[
\boxed{
P(t)=\frac{1000}{1+9e^{-0.4t}}.
}
\]

When does the population reach half the carrying capacity?

Set

\[
P(t)=500.
\]

Then

\[
500=\frac{1000}{1+9e^{-0.4t}}.
\]

So

\[
1+9e^{-0.4t}=2,
\]

and therefore

\[
9e^{-0.4t}=1.
\]

Thus

\[
e^{-0.4t}=\frac19.
\]

Taking logarithms,

\[
-0.4t=-\ln9.
\]

So

\[
\boxed{
t=\frac{\ln9}{0.4}\approx5.49.
}
\]

The population reaches \(500\) after about \(5.49\) time units.

\begin{figure}[h]
\centering
\begin{tikzpicture}
\begin{axis}[
    width=0.84\textwidth,
    height=0.50\textwidth,
    xmin=0, xmax=20,
    ymin=0, ymax=1300,
    axis lines=left,
    xlabel={\(t\)},
    ylabel={\(P(t)\)},
    xtick={0,5,10,15,20},
    ytick={0,500,1000},
    grid=both,
    grid style={line width=.1pt, draw=gray!25},
    legend style={draw=none, at={(0.97,0.35)}, anchor=east}
]
\addplot[very thick, domain=0:20, samples=200] {1000/(1+9*exp(-0.4*x))};
\addlegendentry{\(P_0=100\)}

\addplot[very thick, dashed, domain=0:20, samples=200] {1000/(1-0.1666667*exp(-0.4*x))};
\addlegendentry{\(P_0=1200\)}

\addplot[thick, dotted, domain=0:20, samples=2] {1000};
\addlegendentry{\(P=K\)}

\addplot[only marks, mark=*] coordinates {(5.493,500)};
\node[anchor=west] at (axis cs:5.8,520) {\(P=K/2\)};
\node[anchor=west] at (axis cs:11,1030) {carrying capacity};
\end{axis}
\end{tikzpicture}
\caption{Logistic solutions move toward the carrying capacity.  From below they rise toward \(K\); from above they fall toward \(K\).}
\label{fig:logistic-solutions-ch13}
\end{figure}

The formula gives exact values.  But the most important behavior was already visible
from the differential equation:

\[
0<P<K
\quad\Longrightarrow\quad
P'>0,
\]

and

\[
P>K
\quad\Longrightarrow\quad
P'<0.
\]

So the equation pushes positive solutions toward \(K\).

\subsection{Phase line}
\label{subsec:phase-line}

For an autonomous equation

\[
y'=f(y),
\]

the sign of \(f(y)\) depends only on \(y\).  That means one vertical line can summarize
the motion.

This vertical summary is called a phase line.

For the logistic equation

\[
P'=rP\left(1-\frac{P}{K}\right),
\]

the equilibria occur where

\[
rP\left(1-\frac{P}{K}\right)=0.
\]

Thus

\[
P=0
\qquad\text{and}\qquad
P=K
\]

are equilibrium solutions.

Now read the sign:

\[
P<0
\quad\Longrightarrow\quad
P'<0,
\]

\[
0<P<K
\quad\Longrightarrow\quad
P'>0,
\]

\[
P>K
\quad\Longrightarrow\quad
P'<0.
\]

The phase line is:

\begin{figure}[h]
\centering
\begin{tikzpicture}[xscale=1.25, yscale=1.2, every node/.style={font=\small}]
  \draw[->, thick] (0,-1.0) -- (0,3.0);
  \node[above] at (0,3.0) {\(P\)};

  \draw[thick] (-0.08,0) -- (0.08,0);
  \draw[thick] (-0.08,2) -- (0.08,2);

  \node[left] at (-0.1,0) {\(0\)};
  \node[left] at (-0.1,2) {\(K\)};

  \draw[->, thick] (0,-0.8) -- (0,-0.35);
  \draw[->, thick] (0,0.45) -- (0,1.45);
  \draw[->, thick] (0,2.75) -- (0,2.25);

  \node[right] at (0.15,-0.6) {\(P'<0\)};
  \node[right] at (0.15,0.95) {\(P'>0\)};
  \node[right] at (0.15,2.55) {\(P'<0\)};

  \node[right] at (0.15,0) {unstable};
  \node[right] at (0.15,2) {stable};
\end{tikzpicture}
\caption{The logistic phase line shows how the equation pushes solutions along the \(P\)-axis.}
\label{fig:logistic-phase-line-ch13}
\end{figure}

The arrows tell us the long-term behavior.

If

\[
0<P(0)<K,
\]

then \(P\) rises toward \(K\).

If

\[
P(0)>K,
\]

then \(P\) falls toward \(K\).

If

\[
P(0)=K,
\]

then \(P(t)=K\) forever.

If

\[
P(0)=0,
\]

then \(P(t)=0\) forever.

For a population model, the physically meaningful domain is usually

\[
P\ge0.
\]

The mathematical equation also has behavior for \(P<0\), but negative population is
not part of the model's interpretation.

\begin{quote}
\textbf{Phase line.}
For an autonomous equation

\[
y'=f(y),
\]

a \emph{phase line} marks the equilibrium values \(f(y)=0\) and uses arrows to show
where \(y'>0\) and where \(y'<0\).
\end{quote}

A phase line is a compressed direction field.  Since the slope depends only on \(y\),
the whole plane can be summarized on one vertical axis.

\paragraph{A threshold model.}
The logistic equation says that any small positive population grows.  Some populations
do not behave that way.  If the population is too small, individuals may have trouble
finding mates, protecting each other, or maintaining enough genetic variety.  Below a
threshold, the population may decline.

Let

\[
A=\text{threshold population},
\qquad
K=\text{carrying capacity},
\]

where

\[
0<A<K.
\]

A simple threshold model is

\[
\boxed{
P'=rP\left(1-\frac{P}{K}\right)\left(\frac{P}{A}-1\right),
\qquad r>0.
}
\]

This equation has three equilibria:

\[
P=0,\qquad P=A,\qquad P=K.
\]

The signs are:

\[
0<P<A
\quad\Longrightarrow\quad
P'<0,
\]

\[
A<P<K
\quad\Longrightarrow\quad
P'>0,
\]

\[
P>K
\quad\Longrightarrow\quad
P'<0.
\]

So the phase line is:

\begin{figure}[h]
\centering
\begin{tikzpicture}[xscale=1.25, yscale=1.05, every node/.style={font=\small}]
  \draw[->, thick] (0,-0.5) -- (0,3.3);
  \node[above] at (0,3.3) {\(P\)};

  \draw[thick] (-0.08,0) -- (0.08,0);
  \draw[thick] (-0.08,1.2) -- (0.08,1.2);
  \draw[thick] (-0.08,2.6) -- (0.08,2.6);

  \node[left] at (-0.1,0) {\(0\)};
  \node[left] at (-0.1,1.2) {\(A\)};
  \node[left] at (-0.1,2.6) {\(K\)};

  \draw[->, thick] (0,0.95) -- (0,0.35);
  \draw[->, thick] (0,1.55) -- (0,2.25);
  \draw[->, thick] (0,3.05) -- (0,2.75);

  \node[right] at (0.15,0.65) {\(P'<0\)};
  \node[right] at (0.15,1.9) {\(P'>0\)};
  \node[right] at (0.15,2.95) {\(P'<0\)};

  \node[right] at (0.15,0) {stable from the right};
  \node[right] at (0.15,1.2) {unstable threshold};
  \node[right] at (0.15,2.6) {stable};
\end{tikzpicture}
\caption{In a threshold model, the population must start above \(A\) to move toward the carrying capacity \(K\).}
\label{fig:threshold-phase-line}
\end{figure}

This model has a different story from logistic growth.

If

\[
0<P(0)<A,
\]

the population declines toward extinction.

If

\[
A<P(0)<K,
\]

the population grows toward \(K\).

If

\[
P(0)>K,
\]

the population decreases toward \(K\).

The value \(A\) is a tipping threshold.  Starting just above it and starting just below it
lead to opposite long-term outcomes.

\subsection{Stable and unstable equilibria}
\label{subsec:stable-unstable-equilibria}

An equilibrium solution is a constant solution.  But not all equilibria behave the same
way.

For

\[
y'=f(y),
\]

an equilibrium \(y=c\) occurs when

\[
f(c)=0.
\]

The arrows near \(c\) decide whether nearby solutions move toward \(c\) or away from it.

\begin{quote}
\textbf{Stable and unstable equilibrium.}
An equilibrium \(y=c\) is \emph{stable} if solutions that start near \(c\) move toward
\(c\) as \(t\) increases.

It is \emph{unstable} if solutions that start near \(c\) move away from \(c\) as \(t\)
increases.
\end{quote}

For the logistic equation,

\[
P'=rP\left(1-\frac{P}{K}\right),
\]

the equilibrium \(P=K\) is stable.  Solutions below \(K\) move upward, and solutions
above \(K\) move downward.

The equilibrium \(P=0\) is unstable.  A solution that starts exactly at \(0\) stays there,
but a solution that starts slightly above \(0\) moves away from \(0\).

For the threshold model,

\[
P'=rP\left(1-\frac{P}{K}\right)\left(\frac{P}{A}-1\right),
\]

the equilibrium \(P=A\) is unstable.  A tiny change in the initial condition can change
the long-term outcome.

\[
P(0)=A+\text{small positive amount}
\quad\Longrightarrow\quad
P(t)\to K,
\]

while

\[
P(0)=A-\text{small positive amount}
\quad\Longrightarrow\quad
P(t)\to0.
\]

The threshold \(A\) is not a resting place that attracts the population.  It is a dividing
line.

\begin{quote}
\textbf{What stability means in a model.}
A stable equilibrium is forgiving: small errors in the starting value still lead back to it.

An unstable equilibrium is delicate: small errors in the starting value move the solution
away from it.
\end{quote}

This matters in applications.  A population near a stable carrying capacity can recover
from small disturbances.  A population near an unstable threshold may not.

\subsection{Harvesting and tipping points}
\label{subsec:harvesting-tipping-points}

A fish population may grow logistically, but people may also remove fish by harvesting.
A simple constant-harvest model is

\[
\boxed{
P'=rP\left(1-\frac{P}{K}\right)-h,
}
\]

where

\[
h=\text{harvest rate}.
\]

The term

\[
rP\left(1-\frac{P}{K}\right)
\]

is natural growth.  The term \(h\) is removal.

Equilibria occur when

\[
rP\left(1-\frac{P}{K}\right)=h.
\]

So the equilibrium populations are intersection points between the logistic growth curve
and the horizontal harvest line.

The growth curve

\[
G(P)=rP\left(1-\frac{P}{K}\right)
\]

is a parabola.  Its maximum occurs at

\[
P=\frac K2.
\]

The maximum value is

\[
G\left(\frac K2\right)
=
r\frac K2\left(1-\frac12\right)
=
\boxed{\frac{rK}{4}}.
\]

This number is the largest sustainable constant harvest in the model.

\[
\boxed{
h_{\max}=\frac{rK}{4}.
}
\]

\begin{figure}[h]
\centering
\begin{tikzpicture}
\begin{axis}[
    width=0.82\textwidth,
    height=0.50\textwidth,
    xmin=0, xmax=1000,
    ymin=0, ymax=190,
    axis lines=left,
    xlabel={\(P\)},
    ylabel={rate},
    xtick={0,200,500,800,1000},
    ytick={0,96,150,180},
    grid=both,
    grid style={line width=.1pt, draw=gray!25},
    legend style={draw=none, at={(0.97,0.97)}, anchor=north east}
]
\addplot[very thick, domain=0:1000, samples=200] {0.6*x*(1-x/1000)};
\addlegendentry{\(G(P)=0.6P(1-P/1000)\)}

\addplot[thick, dashed, domain=0:1000, samples=2] {96};
\addlegendentry{\(h=96\)}

\addplot[thick, dotted, domain=0:1000, samples=2] {150};
\addlegendentry{\(h=150\)}

\addplot[thick, dash dot, domain=0:1000, samples=2] {180};
\addlegendentry{\(h=180\)}

\addplot[only marks, mark=*] coordinates {(200,96) (800,96) (500,150)};
\node[anchor=south] at (axis cs:500,152) {maximum sustainable harvest};
\end{axis}
\end{tikzpicture}
\caption{A constant harvest is sustainable only if the harvest line meets the natural growth curve.}
\label{fig:harvesting-growth-curve}
\end{figure}

\paragraph{Example.}
Suppose

\[
r=0.6,
\qquad
K=1000.
\]

Then

\[
h_{\max}=\frac{rK}{4}=\frac{0.6(1000)}{4}=150.
\]

If the harvest rate is

\[
h=96,
\]

then equilibria satisfy

\[
0.6P\left(1-\frac{P}{1000}\right)=96.
\]

Expand:

\[
0.6P-0.0006P^2=96.
\]

Move all terms to one side:

\[
0.0006P^2-0.6P+96=0.
\]

Divide by \(0.0006\):

\[
P^2-1000P+160000=0.
\]

Factor:

\[
(P-200)(P-800)=0.
\]

So the equilibria are

\[
P=200
\qquad\text{and}\qquad
P=800.
\]

To determine stability, examine the sign of

\[
P'=G(P)-96.
\]

When \(P<200\), the harvest is larger than natural growth, so

\[
P'<0.
\]

When \(200<P<800\), natural growth is larger than harvest, so

\[
P'>0.
\]

When \(P>800\), natural growth is again smaller than harvest, so

\[
P'<0.
\]

Thus \(P=200\) is unstable, and \(P=800\) is stable.

The lower equilibrium is a threshold.  If the population falls below \(200\), the model
predicts decline toward extinction.  If the population stays above \(200\), it can recover
toward \(800\).

\paragraph{Three harvesting regimes.}
The model has three cases.

\begin{center}
\begin{tabular}{c|c|c}
\textbf{harvest rate} & \textbf{equilibria} & \textbf{behavior} \\ \hline
\(0<h<\dfrac{rK}{4}\) & two positive equilibria & lower unstable, upper stable \\[8pt]
\(h=\dfrac{rK}{4}\) & one positive equilibrium at \(K/2\) & tipping point \\[8pt]
\(h>\dfrac{rK}{4}\) & no positive equilibrium & harvest exceeds natural growth
\end{tabular}
\end{center}

The middle case is delicate.  If

\[
h=\frac{rK}{4},
\]

then

\[
P'=rP\left(1-\frac{P}{K}\right)-\frac{rK}{4}.
\]

For \(K=1000\) and \(r=0.6\), this becomes

\[
P'=0.6P\left(1-\frac{P}{1000}\right)-150.
\]

The right side is

\[
-0.0006(P-500)^2.
\]

So

\[
P'\le0
\]

for every \(P\), and equality occurs only at \(P=500\).  A population starting above
\(500\) decreases toward \(500\).  A population starting below \(500\) decreases away
from \(500\) toward extinction.

\begin{quote}
\textbf{Model warning.}
The equation

\[
P'=rP\left(1-\frac{P}{K}\right)-h
\]

continues mathematically below \(P=0\), but the population model does not.

If the population reaches extinction, harvesting cannot continue at the same constant
rate.  The physical domain \(P\ge0\) is part of the model.
\end{quote}

The logistic and threshold models show why a differential equation is more than a
formula to solve.  It is a field of tendencies.  It tells what the system tries to do from
each starting value.

The next class of equations has a different structure.  It may include an outside input,
such as salt entering a tank or money deposited into an account.  The equation is not
always separable, but the product rule gives it a way forward.

\section{First-order linear equations}
\label{sec:first-order-linear-equations}

Separable equations can be solved by pulling all \(y\)'s to one side and all \(x\)'s to the
other.  Many useful equations do not separate.

A tank may have salt entering and salt leaving.  A bank account may earn interest while
money is also being deposited.  A cooling object may sit in a room whose temperature is
changing.

These models have a common form:

\[
\text{rate of change}
+
\text{known multiple of the current amount}
=
\text{outside input}.
\]

That is the shape of a first-order linear equation.

\subsection{Standard form}
\label{subsec:first-order-linear-standard-form}

A first-order differential equation is \emph{linear} if the unknown function and its
derivative appear only to the first power and are not multiplied together.

The standard form is

\[
\boxed{
y'+p(t)y=q(t).
}
\]

Here

\[
p(t)
\qquad\text{and}\qquad
q(t)
\]

are known functions of \(t\).  The unknown is \(y(t)\).

\begin{quote}
\textbf{First-order linear equation.}
A first-order linear differential equation has the form

\[
\boxed{
y'+p(t)y=q(t).
}
\]

If \(q(t)=0\), the equation is called \emph{homogeneous}.  If \(q(t)\neq0\), it is called
\emph{nonhomogeneous}.
\end{quote}

The word linear means linear in \(y\) and \(y'\).  The coefficient functions \(p(t)\) and
\(q(t)\) may be complicated functions of \(t\).

\[
y'+(\sin t)y=e^{-t}
\]

is linear.

\[
y'+t^2y=\cos t
\]

is linear.

But

\[
y'=y^2
\]

is not linear, and

\[
y'+y\,y'=t
\]

is not linear, because the unknown function appears nonlinearly.

\paragraph{A first linear example.}
Consider

\[
y'+0.2y=6.
\]

This says that \(y\) is being pushed upward by a constant input \(6\), while the term
\(0.2y\) pulls against the current amount.

We can rewrite it as

\[
y'=6-0.2y.
\]

The equilibrium occurs when

\[
6-0.2y=0.
\]

Thus

\[
y=30.
\]

If \(y<30\), then \(y'>0\), so the solution rises.

If \(y>30\), then \(y'<0\), so the solution falls.

So we expect solutions to approach \(30\).

This particular equation can be solved by shifting the variable:

\[
z=y-30.
\]

Then

\[
z'=y',
\]

and

\[
y'=6-0.2y
=
6-0.2(z+30)
=
6-0.2z-6
=
-0.2z.
\]

So

\[
z'=-0.2z.
\]

Therefore

\[
z=Ce^{-0.2t}.
\]

Since

\[
z=y-30,
\]

we get

\[
\boxed{
y(t)=30+Ce^{-0.2t}.
}
\]

If

\[
y(0)=0,
\]

then

\[
0=30+C,
\]

so

\[
C=-30.
\]

The solution is

\[
\boxed{
y(t)=30\left(1-e^{-0.2t}\right).
}
\]

This example was friendly because the right side was constant.  We need a method that
also works when \(p(t)\) and \(q(t)\) vary with time.

\subsection{Integrating factors}
\label{subsec:integrating-factors}

The product rule says

\[
(\mu y)'=\mu y'+\mu' y.
\]

A linear equation has the left side

\[
y'+p(t)y.
\]

If we multiply the equation by a function \(\mu(t)\), we get

\[
\mu y'+\mu p(t)y=\mu q(t).
\]

We want the left side to be

\[
(\mu y)'.
\]

For that to happen, we need

\[
\mu'=\mu p(t).
\]

That is a separable equation for \(\mu\):

\[
\frac{\mu'}{\mu}=p(t).
\]

Integrate:

\[
\ln|\mu|=\int p(t)\,dt.
\]

So we may choose

\[
\boxed{
\mu(t)=e^{\int p(t)\,dt}.
}
\]

This function is called an integrating factor.

\begin{quote}
\textbf{Integrating factor method.}
To solve

\[
y'+p(t)y=q(t),
\]

choose

\[
\boxed{
\mu(t)=e^{\int p(t)\,dt}.
}
\]

Then

\[
\mu'=\mu p(t),
\]

so multiplying the equation by \(\mu\) gives

\[
(\mu y)'=\mu q(t).
\]

Integrate both sides and solve for \(y\).
\end{quote}

The method is not a trick.  It is the product rule used backward.

\paragraph{Example.}
Solve

\[
y'+2y=e^t,
\qquad
y(0)=1.
\]

Here

\[
p(t)=2,
\qquad
q(t)=e^t.
\]

An integrating factor is

\[
\mu(t)=e^{\int 2\,dt}=e^{2t}.
\]

Multiply the equation by \(e^{2t}\):

\[
e^{2t}y'+2e^{2t}y=e^{3t}.
\]

The left side is

\[
(e^{2t}y)'.
\]

So

\[
(e^{2t}y)'=e^{3t}.
\]

Integrate:

\[
e^{2t}y=\int e^{3t}\,dt=\frac13e^{3t}+C.
\]

Therefore

\[
y=\frac13e^t+Ce^{-2t}.
\]

Use the initial condition:

\[
1=y(0)=\frac13+C.
\]

So

\[
C=\frac23.
\]

Thus

\[
\boxed{
y(t)=\frac13e^t+\frac23e^{-2t}.
}
\]

\paragraph{Check.}
Differentiate:

\[
y'(t)=\frac13e^t-\frac43e^{-2t}.
\]

Then

\[
y'+2y
=
\left(\frac13e^t-\frac43e^{-2t}\right)
+
2\left(\frac13e^t+\frac23e^{-2t}\right).
\]

The exponential decay terms cancel:

\[
-\frac43e^{-2t}+\frac43e^{-2t}=0.
\]

The \(e^t\) terms add:

\[
\frac13e^t+\frac23e^t=e^t.
\]

So

\[
y'+2y=e^t.
\]

Also,

\[
y(0)=\frac13+\frac23=1.
\]

The solution checks.

\subsection{Mixing and finance examples}
\label{subsec:linear-mixing-finance-examples}

Linear equations appear whenever there is a current amount and an outside input.

\paragraph{Example: salt entering and leaving a tank.}
A tank contains \(100\) liters of water.  Brine flows in at \(5\) liters per minute with
concentration

\[
0.2\text{ grams/liter}.
\]

The well-mixed solution flows out at \(5\) liters per minute, so the volume remains
\(100\) liters.

Let

\[
A(t)=\text{grams of salt in the tank after \(t\) minutes}.
\]

Salt enters at the rate

\[
5(0.2)=1
\]

gram per minute.

The concentration in the tank is

\[
\frac{A(t)}{100}
\]

grams per liter, so salt leaves at the rate

\[
5\cdot \frac{A(t)}{100}=\frac{A(t)}{20}
\]

grams per minute.

Thus

\[
A'=1-\frac{A}{20}.
\]

In linear standard form,

\[
\boxed{
A'+\frac1{20}A=1.
}
\]

Use an integrating factor:

\[
\mu(t)=e^{\int (1/20)\,dt}=e^{t/20}.
\]

Multiply by \(e^{t/20}\):

\[
e^{t/20}A'+\frac1{20}e^{t/20}A=e^{t/20}.
\]

The left side is

\[
(e^{t/20}A)'.
\]

So

\[
(e^{t/20}A)'=e^{t/20}.
\]

Integrate:

\[
e^{t/20}A=20e^{t/20}+C.
\]

Therefore

\[
A=20+Ce^{-t/20}.
\]

If the tank initially contains pure water, then

\[
A(0)=0.
\]

So

\[
0=20+C,
\]

and

\[
C=-20.
\]

Thus

\[
\boxed{
A(t)=20\left(1-e^{-t/20}\right).
}
\]

The amount of salt approaches

\[
20\text{ grams}.
\]

That makes sense.  The incoming concentration is \(0.2\) grams per liter, and the tank
holds \(100\) liters.  If the tank eventually matches the incoming concentration, it should
contain

\[
0.2(100)=20
\]

grams of salt.

\begin{figure}[h]
\centering
\begin{tikzpicture}
\begin{axis}[
    width=0.82\textwidth,
    height=0.48\textwidth,
    xmin=0, xmax=100,
    ymin=0, ymax=23,
    axis lines=left,
    xlabel={\(t\) (minutes)},
    ylabel={\(A(t)\) (grams)},
    xtick={0,20,40,60,80,100},
    ytick={0,10,20},
    grid=both,
    grid style={line width=.1pt, draw=gray!25},
]
\addplot[very thick, domain=0:100, samples=200] {20*(1-exp(-x/20))};
\addplot[thick, dotted, domain=0:100, samples=2] {20};
\node[anchor=west] at (axis cs:52,20.4) {equilibrium \(20\) grams};
\node[anchor=west] at (axis cs:24,11) {\(A(t)=20(1-e^{-t/20})\)};
\end{axis}
\end{tikzpicture}
\caption{With a constant incoming concentration, the salt amount approaches an equilibrium.}
\label{fig:linear-mixing-salt}
\end{figure}

\paragraph{Example: continuous deposits with interest.}
A bank account earns continuous interest at annual rate

\[
4\%.
\]

Money is deposited continuously at

\[
1200\text{ dollars per year}.
\]

Let

\[
B(t)=\text{account balance after \(t\) years}.
\]

Interest contributes

\[
0.04B
\]

dollars per year.  Deposits contribute

\[
1200
\]

dollars per year.  Therefore

\[
B'=0.04B+1200.
\]

Write this in standard linear form:

\[
B'-0.04B=1200.
\]

Here

\[
p(t)=-0.04.
\]

An integrating factor is

\[
\mu(t)=e^{\int -0.04\,dt}=e^{-0.04t}.
\]

Multiply:

\[
e^{-0.04t}B'-0.04e^{-0.04t}B=1200e^{-0.04t}.
\]

The left side is

\[
(e^{-0.04t}B)'.
\]

So

\[
(e^{-0.04t}B)'=1200e^{-0.04t}.
\]

Integrate:

\[
e^{-0.04t}B
=
1200\int e^{-0.04t}\,dt
=
1200\left(-\frac1{0.04}e^{-0.04t}\right)+C.
\]

Since

\[
\frac{1200}{0.04}=30000,
\]

we get

\[
e^{-0.04t}B=-30000e^{-0.04t}+C.
\]

Thus

\[
B=-30000+Ce^{0.04t}.
\]

If the account starts empty,

\[
B(0)=0,
\]

then

\[
0=-30000+C,
\]

so

\[
C=30000.
\]

Therefore

\[
\boxed{
B(t)=30000\left(e^{0.04t}-1\right).
}
\]

The formula has a natural interpretation.  The deposits alone would add \(1200t\).  The
interest makes earlier deposits grow, so the balance is larger than \(1200t\) for \(t>0\).

After \(10\) years,

\[
B(10)=30000\left(e^{0.4}-1\right).
\]

Since

\[
e^{0.4}\approx1.4918,
\]

we get

\[
B(10)\approx30000(0.4918)=14754.
\]

The account has about

\[
\boxed{\$14{,}754}
\]

after \(10\) years.

\subsection{Why the integrating factor works}
\label{subsec:why-integrating-factor-works}

The integrating factor method is worth seeing in one clean line.

Start with

\[
y'+p(t)y=q(t).
\]

Suppose we can find a positive function \(\mu(t)\) such that

\[
\mu'=p(t)\mu.
\]

Then multiplying the equation by \(\mu\) gives

\[
\mu y'+\mu p(t)y=\mu q(t).
\]

But

\[
\mu p(t)=\mu',
\]

so the left side becomes

\[
\mu y'+\mu' y.
\]

By the product rule,

\[
\mu y'+\mu' y=(\mu y)'.
\]

Therefore

\[
(\mu y)'=\mu q(t).
\]

Now integrate.

The whole method rests on choosing \(\mu\) so that the left side becomes one derivative.

Since

\[
\frac{\mu'}{\mu}=p(t),
\]

we may take

\[
\mu(t)=e^{\int p(t)\,dt}.
\]

The exponential is useful because its derivative brings down the derivative of the
exponent:

\[
\frac{d}{dt}e^{\int p(t)\,dt}
=
p(t)e^{\int p(t)\,dt}.
\]

So

\[
\mu'=p(t)\mu.
\]

\begin{quote}
\textbf{The product-rule sentence.}
The integrating factor is chosen so that

\[
\mu y'+\mu p(t)y
\]

becomes

\[
(\mu y)'.
\]

That is why the method works.
\end{quote}

There is also a definite-integral version that is useful for initial value problems.

Suppose

\[
y'+p(t)y=q(t),
\qquad
y(t_0)=y_0,
\]

and \(p,q\) are continuous on an interval containing \(t_0\).  Define

\[
\Phi(t)=\int_{t_0}^{t}p(s)\,ds
\]

and

\[
\mu(t)=e^{\Phi(t)}.
\]

Then

\[
\mu(t_0)=e^0=1.
\]

Multiplying the equation by \(\mu\) gives

\[
(\mu y)'=\mu q.
\]

Integrate from \(t_0\) to \(t\):

\[
\mu(t)y(t)-\mu(t_0)y(t_0)
=
\int_{t_0}^{t}\mu(s)q(s)\,ds.
\]

Since \(\mu(t_0)=1\) and \(y(t_0)=y_0\),

\[
\mu(t)y(t)-y_0
=
\int_{t_0}^{t}\mu(s)q(s)\,ds.
\]

Therefore

\[
\boxed{
y(t)
=
e^{-\Phi(t)}
\left[
y_0+\int_{t_0}^{t}e^{\Phi(s)}q(s)\,ds
\right].
}
\]

This formula is not always the easiest way to compute, but it tells the whole structure:

\[
\text{initial amount}
+
\text{accumulated weighted input}.
\]

\paragraph{Why the interval matters.}
The formula assumes that \(p(t)\) and \(q(t)\) are continuous on the interval where we
solve the equation.  If a coefficient breaks, the solution may have to be handled on
separate intervals.

Consider

\[
y'+\frac1t y=1.
\]

This is linear, but the coefficient

\[
\frac1t
\]

is not defined at

\[
t=0.
\]

So we solve on intervals that avoid \(0\), such as

\[
(0,\infty)
\]

or

\[
(-\infty,0).
\]

On \(t>0\), an integrating factor is

\[
\mu(t)=e^{\int (1/t)\,dt}=e^{\ln t}=t.
\]

Then

\[
ty'+y=t,
\]

so

\[
(ty)'=t.
\]

Integrate:

\[
ty=\frac{t^2}{2}+C.
\]

Thus

\[
\boxed{
y=\frac t2+\frac{C}{t},
\qquad t>0.
}
\]

The formula cannot be continued through \(t=0\) because of the term \(C/t\), and the
original equation itself is not defined at \(0\).

\begin{quote}
\textbf{Coefficient warning.}
A solution of a differential equation lives on an interval where the equation makes
sense.  Singular coefficients split the real line into separate problems.
\end{quote}

First-order equations already show several personalities.  Some are separable.  Some
are linear.  Some are best understood by a phase line.  Some require numerical methods.

But many laws of motion do not stop at the first derivative.  A spring, a pendulum in a
small-angle approximation, and a vibrating string involve acceleration.  Their laws begin
with

\[
y''.
\]

That is the next door: differential equations where the unknown function is governed by
its second derivative.
\section{Second-order equations and vibration}
\label{sec:second-order-equations-vibration}

The first-order equations in this chapter described quantities whose rates were governed
by their current state or by an outside input.  But many physical laws do not begin with
velocity.  They begin with acceleration.

A spring does not merely say where a mass is.  It pulls harder when the mass is farther
from equilibrium.  Newton's second law then connects that force to acceleration.

\[
\text{force}
\quad\Longrightarrow\quad
\text{acceleration}
\quad\Longrightarrow\quad
\text{second derivative}.
\]

That is why vibrating systems naturally lead to second-order differential equations.

\subsection{Simple harmonic motion}
\label{subsec:simple-harmonic-motion}

Imagine a mass attached to a spring on a frictionless table.  Pull the mass to the right,
let go, and it moves back and forth.

Let

\[
y(t)=\text{displacement from equilibrium at time }t.
\]

So

\[
y=0
\]

means the mass is at its resting position.  Positive \(y\) means it is displaced to the
right.  Negative \(y\) means it is displaced to the left.

A spring obeying Hooke's law pulls back toward equilibrium:

\[
F=-ky,
\qquad k>0.
\]

The constant \(k\) measures the stiffness of the spring.  The minus sign is the important
part.  If \(y>0\), the force is negative.  If \(y<0\), the force is positive.  The force always
points back toward \(0\).

Newton's second law says

\[
F=ma,
\]

where

\[
a=y''.
\]

Therefore

\[
my''=-ky.
\]

Rearrange:

\[
y''+\frac{k}{m}y=0.
\]

It is customary to write

\[
\omega^2=\frac{k}{m},
\]

where

\[
\omega>0.
\]

Then the spring equation becomes

\[
\boxed{
y''+\omega^2 y=0.
}
\]

This equation says:

\[
\boxed{
\text{acceleration is proportional to displacement, but in the opposite direction.}
}
\]

That is the mathematical heart of simple harmonic motion.

\paragraph{Example.}
Check that

\[
y(t)=3\cos(2t)
\]

solves

\[
y''+4y=0.
\]

Differentiate:

\[
y'(t)=-6\sin(2t),
\]

and

\[
y''(t)=-12\cos(2t).
\]

Now compute the left side:

\[
y''+4y
=
-12\cos(2t)+4\bigl(3\cos(2t)\bigr)
=
0.
\]

So

\[
\boxed{
y(t)=3\cos(2t)
}
\]

is a solution.

The initial position is

\[
y(0)=3,
\]

and the initial velocity is

\[
y'(0)=0.
\]

That matches the picture of a mass pulled \(3\) units from equilibrium and released
from rest.

\begin{figure}[h]
\centering
\begin{tikzpicture}
\begin{axis}[
    width=0.84\textwidth,
    height=0.48\textwidth,
    xmin=0, xmax=7,
    ymin=-3.6, ymax=3.6,
    axis lines=middle,
    xlabel={\(t\)},
    ylabel={\(y(t)\)},
    xtick={0,1.5708,3.1416,4.7124,6.2832},
    xticklabels={\(0\),\(\frac{\pi}{2}\),\(\pi\),\(\frac{3\pi}{2}\),\(2\pi\)},
    ytick={-3,0,3},
    grid=both,
    grid style={line width=.1pt, draw=gray!25},
]
\addplot[very thick, domain=0:7, samples=300] {3*cos(deg(2*x))};
\addplot[thick, dotted, domain=0:7, samples=2] {0};
\node[anchor=south west] at (axis cs:0.15,3) {\(y=3\cos(2t)\)};
\node[anchor=west] at (axis cs:3.35,0.35) {equilibrium};
\end{axis}
\end{tikzpicture}
\caption{Simple harmonic motion oscillates around the equilibrium \(y=0\).}
\label{fig:simple-harmonic-motion-cos}
\end{figure}

The graph repeats.  The mass returns to the same position and velocity again and again.

The number \(\omega\) is called the angular frequency.  The period is

\[
\boxed{
T=\frac{2\pi}{\omega}.
}
\]

For \(y=3\cos(2t)\), we have \(\omega=2\), so

\[
T=\frac{2\pi}{2}=\pi.
\]

\subsection{\texorpdfstring{The equation \(y''+\omega^2y=0\)}{The equation y'' + omega^2 y = 0}}
\label{subsec:equation-ypp-omega2-y}

The equation

\[
y''+\omega^2y=0
\]

is friendly because sine and cosine repeat under differentiation.

\[
(\cos \omega t)'=-\omega\sin\omega t,
\]

\[
(\cos \omega t)''=-\omega^2\cos\omega t.
\]

So

\[
y=\cos\omega t
\]

satisfies

\[
y''+\omega^2y=0.
\]

The same is true for

\[
y=\sin\omega t,
\]

because

\[
(\sin \omega t)''=-\omega^2\sin\omega t.
\]

Since the equation is linear, any linear combination is also a solution:

\[
\boxed{
y(t)=C_1\cos(\omega t)+C_2\sin(\omega t).
}
\]

This is the general solution.

\begin{quote}
\textbf{Simple harmonic oscillator.}
The differential equation

\[
\boxed{
y''+\omega^2y=0,
\qquad \omega>0,
}
\]

has general solution

\[
\boxed{
y(t)=C_1\cos(\omega t)+C_2\sin(\omega t).
}
\]
\end{quote}

The constants are determined by two initial conditions.  A second-order equation needs
two starting values because motion needs both position and velocity.

Suppose

\[
y(0)=y_0
\qquad\text{and}\qquad
y'(0)=v_0.
\]

From

\[
y(t)=C_1\cos(\omega t)+C_2\sin(\omega t),
\]

we get

\[
y(0)=C_1.
\]

So

\[
C_1=y_0.
\]

Differentiate:

\[
y'(t)=-\omega C_1\sin(\omega t)+\omega C_2\cos(\omega t).
\]

Then

\[
y'(0)=\omega C_2.
\]

So

\[
C_2=\frac{v_0}{\omega}.
\]

Thus the solution of

\[
y''+\omega^2y=0,
\qquad
y(0)=y_0,
\qquad
y'(0)=v_0
\]

is

\[
\boxed{
y(t)=y_0\cos(\omega t)+\frac{v_0}{\omega}\sin(\omega t).
}
\]

\paragraph{Example.}
Solve

\[
y''+9y=0,
\qquad
y(0)=2,
\qquad
y'(0)=6.
\]

Here

\[
\omega^2=9,
\qquad
\omega=3.
\]

The solution is

\[
y(t)=C_1\cos(3t)+C_2\sin(3t).
\]

The initial position gives

\[
C_1=2.
\]

Differentiate:

\[
y'(t)=-3C_1\sin(3t)+3C_2\cos(3t).
\]

Then

\[
y'(0)=3C_2.
\]

Since \(y'(0)=6\),

\[
3C_2=6,
\]

so

\[
C_2=2.
\]

Therefore

\[
\boxed{
y(t)=2\cos(3t)+2\sin(3t).
}
\]

\paragraph{Energy check.}
The undamped spring keeps oscillating because no energy is lost.

For

\[
my''+ky=0,
\]

define the energy

\[
E(t)=\frac12m\bigl(y'(t)\bigr)^2+\frac12k\bigl(y(t)\bigr)^2.
\]

The first term is kinetic energy.  The second is spring potential energy.

Differentiate:

\[
E'(t)
=
m y'y''+k yy'.
\]

Factor out \(y'\):

\[
E'(t)=y'(my''+ky).
\]

But the equation says

\[
my''+ky=0.
\]

Therefore

\[
E'(t)=0.
\]

So

\[
\boxed{
E(t)\text{ is constant.}
}
\]

The energy moves back and forth between motion and spring stretch.  It does not
disappear.

\subsection{Damping}
\label{subsec:damping}

Real vibrations usually die down.  A guitar string grows quiet.  A shock absorber settles
after a bump.  A swinging door stops swinging.

The missing ingredient is damping.

Damping is a force that opposes velocity.  A simple model is

\[
F_{\text{damping}}=-c y',
\qquad c>0.
\]

The spring force is still

\[
F_{\text{spring}}=-ky.
\]

Newton's law gives

\[
my''=-cy'-ky.
\]

So

\[
\boxed{
my''+cy'+ky=0.
}
\]

This is a second-order linear differential equation with constant coefficients.  The new
term

\[
cy'
\]

is responsible for energy loss.

\begin{quote}
\textbf{Damped oscillator.}
A damped spring-mass system is modeled by

\[
\boxed{
my''+cy'+ky=0,
}
\]

where

\[
m>0,\qquad c>0,\qquad k>0.
\]

The mass \(m\) resists acceleration, the damping coefficient \(c\) resists velocity, and
the spring constant \(k\) resists displacement.
\end{quote}

For light damping, the motion still oscillates, but the oscillations shrink.  One common
form of such a solution is

\[
y(t)=e^{-\beta t}\bigl(A\cos(\mu t)+B\sin(\mu t)\bigr),
\]

where

\[
\beta>0.
\]

The exponential factor

\[
e^{-\beta t}
\]

is the shrinking envelope.  The sine and cosine still oscillate, but their amplitude decays.

\paragraph{Example.}
Check that

\[
y(t)=e^{-0.2t}\cos(2t)
\]

solves

\[
y''+0.4y'+4.04y=0.
\]

First compute

\[
y'(t)=e^{-0.2t}\bigl(-0.2\cos(2t)-2\sin(2t)\bigr).
\]

Differentiate again:

\[
y''(t)=e^{-0.2t}\bigl(-3.96\cos(2t)+0.8\sin(2t)\bigr).
\]

Now substitute:

\[
y''+0.4y'+4.04y
\]

equals

\[
e^{-0.2t}\bigl(-3.96\cos(2t)+0.8\sin(2t)\bigr)
\]

\[
+0.4e^{-0.2t}\bigl(-0.2\cos(2t)-2\sin(2t)\bigr)
+
4.04e^{-0.2t}\cos(2t).
\]

Collect the cosine terms:

\[
-3.96-0.08+4.04=0.
\]

Collect the sine terms:

\[
0.8-0.8=0.
\]

So the expression is \(0\).  Therefore the function solves the damped oscillator
equation.

\begin{figure}[h]
\centering
\begin{tikzpicture}
\begin{axis}[
    width=0.84\textwidth,
    height=0.48\textwidth,
    xmin=0, xmax=18,
    ymin=-1.2, ymax=1.2,
    axis lines=middle,
    xlabel={\(t\)},
    ylabel={\(y(t)\)},
    xtick={0,5,10,15},
    ytick={-1,0,1},
    grid=both,
    grid style={line width=.1pt, draw=gray!25},
    legend style={draw=none, at={(0.98,0.97)}, anchor=north east}
]
\addplot[very thick, domain=0:18, samples=500] {exp(-0.2*x)*cos(deg(2*x))};
\addlegendentry{\(e^{-0.2t}\cos(2t)\)}
\addplot[thick, dashed, domain=0:18, samples=200] {exp(-0.2*x)};
\addlegendentry{\(e^{-0.2t}\)}
\addplot[thick, dashed, domain=0:18, samples=200] {-exp(-0.2*x)};
\end{axis}
\end{tikzpicture}
\caption{Damping multiplies the oscillation by a decaying envelope.}
\label{fig:damped-oscillation}
\end{figure}

The solution still crosses the equilibrium many times, but each swing is smaller than the
previous one.

\paragraph{Energy with damping.}
For the damped equation

\[
my''+cy'+ky=0,
\]

use the same energy

\[
E(t)=\frac12m(y')^2+\frac12ky^2.
\]

Then

\[
E'(t)=y'(my''+ky).
\]

From the differential equation,

\[
my''+ky=-cy'.
\]

Therefore

\[
E'(t)=y'(-cy')=-c(y')^2.
\]

Since

\[
c>0,
\]

we have

\[
\boxed{
E'(t)\le0.
}
\]

Energy decreases, except at instants when \(y'=0\).  The damping term drains energy
from the motion.

\subsection{Driven oscillation and resonance}
\label{subsec:driven-oscillation-resonance}

A system can also be forced from outside.  A child on a swing is pushed.  A bridge is hit
by periodic wind.  A guitar string is driven by a vibrating body.

The simplest driven oscillator has the form

\[
\boxed{
y''+\omega^2y=A\cos(\gamma t).
}
\]

The left side describes the natural oscillator.  The right side is an outside periodic force.

The number

\[
\omega
\]

is the natural angular frequency.  The number

\[
\gamma
\]

is the driving angular frequency.

If

\[
\gamma\neq\omega,
\]

try a particular solution of the form

\[
y_p(t)=C\cos(\gamma t).
\]

Then

\[
y_p''(t)=-C\gamma^2\cos(\gamma t).
\]

Substitute into the equation:

\[
y_p''+\omega^2y_p
=
C(\omega^2-\gamma^2)\cos(\gamma t).
\]

We want this to equal

\[
A\cos(\gamma t).
\]

So

\[
C(\omega^2-\gamma^2)=A,
\]

and therefore

\[
\boxed{
C=\frac{A}{\omega^2-\gamma^2}.
}
\]

Thus, when \(\gamma\neq\omega\),

\[
\boxed{
y_p(t)=\frac{A}{\omega^2-\gamma^2}\cos(\gamma t).
}
\]

The denominator tells the story.  If the driving frequency \(\gamma\) is close to the
natural frequency \(\omega\), then

\[
|\omega^2-\gamma^2|
\]

is small, and the response amplitude can be large.

\begin{quote}
\textbf{Resonance, first view.}
A system resonates when the driving frequency is close to the natural frequency.  The
forcing then reinforces the natural motion instead of fighting it.
\end{quote}

At exact resonance, the trial form above fails because the denominator becomes \(0\).
For

\[
y''+\omega^2y=A\cos(\omega t),
\]

one particular solution is

\[
\boxed{
y_p(t)=\frac{A}{2\omega}t\sin(\omega t).
}
\]

The factor \(t\) means that the amplitude grows with time in the ideal undamped model.

\begin{figure}[h]
\centering
\begin{tikzpicture}
\begin{axis}[
    width=0.82\textwidth,
    height=0.48\textwidth,
    xmin=0, xmax=3.5,
    ymin=0, ymax=6.2,
    axis lines=left,
    xlabel={driving frequency \(\gamma\)},
    ylabel={response size},
    xtick={0,1,2,3},
    ytick={0,2,4,6},
    grid=both,
    grid style={line width=.1pt, draw=gray!25},
]
\addplot[very thick, domain=0:1.88, samples=200] {1/(4-x^2)};
\addplot[very thick, domain=2.12:3.45, samples=200] {1/(x^2-4)};
\draw[dashed] (axis cs:2,0) -- (axis cs:2,6);
\node[anchor=south] at (axis cs:2,6) {\(\gamma=\omega\)};
\node[anchor=west] at (axis cs:2.15,4.4) {large response};
\node[anchor=west] at (axis cs:0.3,0.65) {\(\omega=2,\ A=1\)};
\end{axis}
\end{tikzpicture}
\caption{In the undamped model \(y''+\omega^2y=A\cos(\gamma t)\), the response grows large as \(\gamma\) approaches \(\omega\).}
\label{fig:resonance-response}
\end{figure}

Real systems have damping, so the amplitude does not grow forever.  But the warning
survives: forcing near the natural frequency can create a much larger response than
forcing far from it.

\begin{quote}
\textbf{Model warning.}
The equation

\[
y''+\omega^2y=A\cos(\omega t)
\]

predicts unbounded growth at resonance because it ignores damping and structural
limits.  It is a warning model, not a complete physical description.
\end{quote}

Second-order equations already show that the state of a system may need more than one
number.  A spring is not determined by position alone.  To know what happens next, we
also need velocity.

That idea leads naturally to systems of differential equations.

\section{Systems and interaction models}
\label{sec:systems-interaction-models}

A first-order equation such as

\[
y'=f(y)
\]

uses one number to describe the current state.  A spring needs two numbers: position and
velocity.  Two interacting populations need two numbers: one for each population.

When the state has several components, one differential equation becomes a system.

\[
\text{one changing quantity}
\quad\Longrightarrow\quad
\text{one differential equation},
\]

\[
\text{several changing quantities}
\quad\Longrightarrow\quad
\text{several linked differential equations}.
\]

\subsection{Two populations}
\label{subsec:two-populations}

A single population model may look like

\[
P'=rP.
\]

That can describe a species living alone with plentiful resources.  But species do not
usually live alone.  They compete, cooperate, eat, flee, and depend on one another.

Let

\[
R(t)=\text{number of rabbits at time }t,
\]

and

\[
W(t)=\text{number of wolves at time }t.
\]

The rabbits are prey.  The wolves are predators.

If there were no wolves, the rabbits might grow approximately exponentially:

\[
R'=0.8R.
\]

If there were no rabbits, the wolves might decline:

\[
W'=-0.6W.
\]

But when both are present, they meet.  The number of encounters should increase when
there are more rabbits and also when there are more wolves.  A simple model makes the
encounter term proportional to

\[
RW.
\]

Predation decreases the rabbit population and increases the wolf population.  One model
is

\[
\boxed{
\begin{aligned}
R'&=0.8R-0.02RW,\\
W'&=-0.6W+0.01RW.
\end{aligned}
}
\]

The first equation says that rabbits reproduce, but are eaten.  The second says that wolves
decline without food, but benefit from encounters with rabbits.

This is not one equation.  It is a system.

\begin{quote}
\textbf{System of differential equations.}
A \emph{system of differential equations} is a set of differential equations for several
unknown functions.

A solution of

\[
\begin{aligned}
R'&=F(R,W),\\
W'&=G(R,W)
\end{aligned}
\]

is a pair of functions

\[
R(t),\qquad W(t)
\]

that satisfies both equations on the same time interval.
\end{quote}

The equations are coupled.  The derivative \(R'\) depends on \(W\), and the derivative
\(W'\) depends on \(R\).  We cannot solve one equation first and then the other.  The two
populations move together.

\paragraph{Equilibria.}
An equilibrium occurs when both derivatives are zero:

\[
R'=0
\qquad\text{and}\qquad
W'=0.
\]

For the model

\[
R'=0.8R-0.02RW=R(0.8-0.02W),
\]

so

\[
R'=0
\]

when

\[
R=0
\qquad\text{or}\qquad
W=40.
\]

Also,

\[
W'=-0.6W+0.01RW=W(-0.6+0.01R),
\]

so

\[
W'=0
\]

when

\[
W=0
\qquad\text{or}\qquad
R=60.
\]

Thus there are two biologically meaningful equilibria:

\[
(R,W)=(0,0)
\]

and

\[
(R,W)=(60,40).
\]

The equilibrium

\[
(0,0)
\]

means no rabbits and no wolves.  The equilibrium

\[
(60,40)
\]

means the two populations balance each other in the model.

\subsection{Predator-prey equations as a picture}
\label{subsec:predator-prey-picture}

For one autonomous equation,

\[
y'=f(y),
\]

a phase line summarized the motion.  For two autonomous equations,

\[
\begin{aligned}
R'&=F(R,W),\\
W'&=G(R,W),
\end{aligned}
\]

the state is a point in the \((R,W)\)-plane.  The differential equation assigns a velocity
vector to each point:

\[
(R,W)
\quad\longmapsto\quad
(R',W').
\]

So the phase line becomes a phase plane.

For the predator-prey model,

\[
\begin{aligned}
R'&=R(0.8-0.02W),\\
W'&=W(-0.6+0.01R),
\end{aligned}
\]

the lines

\[
W=40
\]

and

\[
R=60
\]

divide the plane into four regions.

If

\[
W<40,
\]

then

\[
R'>0.
\]

If

\[
W>40,
\]

then

\[
R'<0.
\]

If

\[
R<60,
\]

then

\[
W'<0.
\]

If

\[
R>60,
\]

then

\[
W'>0.
\]

The signs give the direction of motion.

\[
\begin{array}{c|c|c|c}
\textbf{region} & \textbf{sign of \(R'\)} & \textbf{sign of \(W'\)} & \textbf{motion} \\ \hline
R<60,\ W<40 & + & - & \text{right and down} \\
R>60,\ W<40 & + & + & \text{right and up} \\
R>60,\ W>40 & - & + & \text{left and up} \\
R<60,\ W>40 & - & - & \text{left and down}
\end{array}
\]

\begin{figure}[h]
\centering
\begin{tikzpicture}
\begin{axis}[
    width=0.82\textwidth,
    height=0.56\textwidth,
    xmin=0, xmax=120,
    ymin=0, ymax=85,
    axis lines=left,
    xlabel={rabbits \(R\)},
    ylabel={wolves \(W\)},
    xtick={0,60,120},
    ytick={0,40,80},
    grid=both,
    grid style={line width=.1pt, draw=gray!25},
]
% nullclines
\addplot[thick, dashed, domain=0:120, samples=2] {40};
\draw[thick, dashed] (axis cs:60,0) -- (axis cs:60,85);

\node[anchor=south west] at (axis cs:63,40) {equilibrium \((60,40)\)};
\addplot[only marks, mark=*] coordinates {(60,40)};

% approximate closed orbit
\addplot[very thick, domain=0:360, samples=240]
  ({60+35*cos(x)}, {40+22*sin(x)});

% direction arrows
\draw[->, thick] (axis cs:25,20) -- (axis cs:38,12);
\draw[->, thick] (axis cs:90,20) -- (axis cs:104,31);
\draw[->, thick] (axis cs:95,65) -- (axis cs:80,75);
\draw[->, thick] (axis cs:25,65) -- (axis cs:12,54);

\node[anchor=west] at (axis cs:63,7) {\(R=60\)};
\node[anchor=west] at (axis cs:5,42) {\(W=40\)};
\node[anchor=west] at (axis cs:72,18) {prey abundant};
\node[anchor=west] at (axis cs:72,12) {predators rise};
\end{axis}
\end{tikzpicture}
\caption{A predator-prey system is read in the phase plane.  The point \((R(t),W(t))\) moves according to the vector \((R',W')\).}
\label{fig:predator-prey-phase-plane}
\end{figure}

The picture suggests a cycle.

When wolves are few, rabbits increase.  Then the abundance of rabbits lets wolves
increase.  More wolves reduce the rabbits.  With fewer rabbits, wolves decline.  Then the
rabbits can recover.

The equations do not describe one population rising or falling by itself.  They describe
a moving point in the plane.

\paragraph{A numerical glimpse.}
Suppose

\[
R(0)=80,
\qquad
W(0)=20.
\]

Then

\[
R'(0)=80(0.8-0.02(20))=80(0.4)=32,
\]

and

\[
W'(0)=20(-0.6+0.01(80))=20(0.2)=4.
\]

So at the starting point,

\[
(R',W')=(32,4).
\]

Both populations are increasing, but rabbits are increasing faster.  In the phase plane,
the point initially moves mostly to the right and slightly upward.

At a different point, say

\[
R=30,\qquad W=70,
\]

we get

\[
R'=30(0.8-0.02(70))=30(-0.6)=-18,
\]

and

\[
W'=70(-0.6+0.01(30))=70(-0.3)=-21.
\]

Both populations are decreasing.  The point moves left and down.

This is the system version of a direction field.  Instead of little line segments with a
slope, we draw little arrows with two components.

\begin{quote}
\textbf{Phase plane.}
For a system

\[
\begin{aligned}
x'&=F(x,y),\\
y'&=G(x,y),
\end{aligned}
\]

the \emph{phase plane} is the \((x,y)\)-plane together with direction arrows

\[
(F(x,y),G(x,y)).
\]

A solution traces a curve in this plane as time passes.
\end{quote}

The phase plane does not usually give exact formulas for \(R(t)\) and \(W(t)\).  It gives
the shape of the interaction.

That is often the first thing we need.

\subsection{Why systems lead naturally to multivariable calculus}
\label{subsec:systems-lead-to-multivariable-calculus}

A system changes the meaning of state.

For one equation,

\[
y'=f(y),
\]

the state is one number:

\[
y.
\]

For two equations,

\[
\begin{aligned}
x'&=F(x,y),\\
y'&=G(x,y),
\end{aligned}
\]

the state is a point:

\[
(x,y).
\]

The right side is not a function of one variable.  It is built from functions of two
variables:

\[
F(x,y)
\qquad\text{and}\qquad
G(x,y).
\]

That is the doorway to multivariable calculus.

Even a second-order equation has this structure hiding inside it.  Consider

\[
y''+\omega^2y=0.
\]

Let

\[
v=y'.
\]

Then

\[
y'=v.
\]

Also, from the differential equation,

\[
v'=y''=-\omega^2y.
\]

So the second-order equation is equivalent to the first-order system

\[
\boxed{
\begin{aligned}
y'&=v,\\
v'&=-\omega^2y.
\end{aligned}
}
\]

The state is now

\[
(y,v),
\]

position and velocity together.

For \(\omega=1\), the system is

\[
\begin{aligned}
y'&=v,\\
v'&=-y.
\end{aligned}
\]

The curves in the \((y,v)\)-plane are circles, because the quantity

\[
y^2+v^2
\]

stays constant.  Indeed,

\[
\frac{d}{dt}(y^2+v^2)=2yy'+2vv'.
\]

Substitute

\[
y'=v,
\qquad
v'=-y:
\]

\[
2y v+2v(-y)=0.
\]

So

\[
y^2+v^2=\text{constant}.
\]

The oscillator moves around closed curves in the position-velocity plane.

\begin{figure}[h]
\centering
\begin{tikzpicture}
\begin{axis}[
    width=0.62\textwidth,
    height=0.62\textwidth,
    xmin=-3.2, xmax=3.2,
    ymin=-3.2, ymax=3.2,
    axis lines=middle,
    xlabel={position \(y\)},
    ylabel={velocity \(v\)},
    xtick={-3,-2,-1,0,1,2,3},
    ytick={-3,-2,-1,0,1,2,3},
    grid=both,
    grid style={line width=.1pt, draw=gray!25},
]
% vector field for y'=v, v'=-y
\foreach \xx in {-3,-2,-1,0,1,2,3} {
  \foreach \yy in {-3,-2,-1,0,1,2,3} {
    \pgfmathsetmacro{\vx}{\yy}
    \pgfmathsetmacro{\vy}{-\xx}
    \pgfmathsetmacro{\s}{sqrt(\vx*\vx+\vy*\vy)}
    \ifdim \s pt > 0pt
      \pgfmathsetmacro{\dx}{0.22*\vx/\s}
      \pgfmathsetmacro{\dy}{0.22*\vy/\s}
      
      % Pre-calculate target coordinates
      \pgfmathsetmacro{\xTwo}{\xx+\dx}
      \pgfmathsetmacro{\yTwo}{\yy+\dy}
      
      % Expand variables into hard numbers before plotting
      \edef\temp{%
        \noexpand\draw[gray!65,->] (axis cs:\xx,\yy) -- (axis cs:\xTwo,\yTwo);%
      }
      % Execute the expanded command
      \temp
    \fi
  }
}
\addplot[very thick, domain=0:360, samples=240] ({2*cos(x)}, {2*sin(x)});
\addplot[very thick, dashed, domain=0:360, samples=240] ({1.2*cos(x)}, {1.2*sin(x)});
\node[anchor=west] at (axis cs:1.6,1.6) {\(y^2+v^2=\text{constant}\)};
\end{axis}
\end{tikzpicture}
\caption{The simple oscillator becomes a first-order system in the position-velocity plane.}
\label{fig:oscillator-phase-plane}
\end{figure}

This is a different kind of graph from \(y(t)\) versus \(t\).  It does not show time on an
axis.  It shows the state of the system.  Time is the parameter that traces the curve.

\begin{quote}
\textbf{The bridge.}
A solution of a system is a moving point:

\[
t\longmapsto (x(t),y(t)).
\]

So differential equations lead naturally to parametrized curves, vector fields, and
functions of several variables.
\end{quote}

This is where single-variable calculus begins to lean toward the next subject.  The
derivative of one function was enough for motion on a line.  Motion in a plane needs two
coordinate functions.  Interactions need functions of several variables.  The geometry of
change begins to expand.

Before leaving differential equations, we should collect the warnings that protect the
methods we have used: solutions can be lost by division, numerical methods can invent
false behavior, and the same differential equation can describe many different histories
until an initial condition chooses one.
\section{Warning examples}
\label{sec:de-warning-examples}

Differential equations are laws written as derivatives.  In this chapter, we learned to
read those laws in several ways:

\[
\text{direction fields},
\qquad
\text{Euler steps},
\qquad
\text{separation},
\qquad
\text{integrating factors},
\qquad
\text{phase lines}.
\]

Each method is useful.  Each method also has a danger.

Before leaving differential equations, we collect three warnings.  They are not side
remarks.  They protect the meaning of the whole chapter.

\[
\boxed{
\text{Algebra can lose solutions.}
}
\]

\[
\boxed{
\text{Numerical methods can create false behavior.}
}
\]

\[
\boxed{
\text{A differential equation is not complete without enough starting information.}
}
\]

\subsection{A solution lost by dividing by zero}
\label{subsec:solution-lost-by-dividing-zero}

Separation often begins by dividing by an expression involving the unknown function.
That division may be illegal at some values of the solution.  Those values often give
constant solutions.

Start with a simple equation:

\[
y'=y^2.
\]

If \(y\neq0\), we can separate:

\[
\frac{dy}{y^2}=dt.
\]

Integrate:

\[
\int y^{-2}\,dy=\int 1\,dt.
\]

Thus

\[
-\frac1y=t+C.
\]

Solving for \(y\), we get

\[
y=-\frac{1}{t+C}.
\]

Equivalently, after renaming the constant,

\[
\boxed{
y=\frac{1}{C-t}.
}
\]

These are nonzero solutions.  For example, if

\[
y(0)=1,
\]

then

\[
1=\frac1C,
\]

so \(C=1\), and

\[
\boxed{
y(t)=\frac{1}{1-t}.
}
\]

This solution increases and becomes unbounded as \(t\to1^{-}\).  The differential
equation predicts faster growth when \(y\) is larger, and that feedback becomes
violent.

But the separation step divided by \(y^2\).  It excluded the case

\[
y=0.
\]

Check that case separately.  The constant function

\[
y(t)=0
\]

has derivative

\[
y'(t)=0,
\]

and the right side is

\[
y(t)^2=0.
\]

So

\[
\boxed{
y(t)=0
}
\]

is also a solution of

\[
y'=y^2.
\]

It was lost by dividing by \(y^2\).

\begin{figure}[h]
\centering
\begin{tikzpicture}
\begin{axis}[
    width=0.82\textwidth,
    height=0.50\textwidth,
    xmin=-2.0, xmax=2.0,
    ymin=-3.2, ymax=3.2,
    axis lines=middle,
    xlabel={\(t\)},
    ylabel={\(y\)},
    xtick={-2,-1,0,1,2},
    ytick={-3,-2,-1,0,1,2,3},
    grid=both,
    grid style={line width=.1pt, draw=gray!25},
    legend style={draw=none, at={(0.98,0.98)}, anchor=north east}
]
\addplot[very thick, domain=-2:0.85, samples=200] {1/(1-x)};
\addlegendentry{\(y(0)=1\)}
\addplot[very thick, dashed, domain=-0.85:2, samples=200] {1/(-1-x)};
\addlegendentry{\(y(0)=-1\)}
\addplot[thick, dotted, domain=-2:2, samples=2] {0};
\addlegendentry{\(y=0\)}
\draw[dashed] (axis cs:1,-3.2) -- (axis cs:1,3.2);
\node[anchor=west] at (axis cs:1.05,2.4) {\(t=1\)};
\end{axis}
\end{tikzpicture}
\caption{The separated family \(y=1/(C-t)\) misses the equilibrium solution \(y=0\).}
\label{fig:yprime-y-squared-lost-solution}
\end{figure}

The same warning appears in population models.  For example,

\[
P'=rP\left(1-\frac{P}{K}\right)
\]

has equilibrium solutions

\[
P=0
\qquad\text{and}\qquad
P=K.
\]

If we divide by

\[
P\left(1-\frac{P}{K}\right),
\]

then we have already assumed

\[
P\neq0
\qquad\text{and}\qquad
P\neq K.
\]

Those two constant solutions must be checked separately.

\begin{quote}
\textbf{Division rule for separable equations.}
When separating variables, if you divide by an expression involving \(y\), first ask:

\[
\text{For which constant values of \(y\) does this expression equal zero?}
\]

Those values are candidates for equilibrium solutions.
\end{quote}

This is not only a technical warning.  In models, equilibrium solutions often carry the
main interpretation.  Extinction, carrying capacity, room temperature, and chemical
balance all appear as constant solutions.

Losing them would lose the story.

\subsection{Euler's method follows the wrong behavior with a large step}
\label{subsec:euler-wrong-behavior-large-step}

Euler's method follows local slopes one step at a time:

\[
y_{k+1}=y_k+hF(x_k,y_k).
\]

The method is built from local linear approximation.  That is its strength.  It is also its
weakness.  If the step is too large, the tangent-line instruction may no longer represent
the curve.

Consider the initial value problem

\[
y'=-4y,
\qquad
y(0)=1.
\]

The exact solution is

\[
\boxed{
y=e^{-4x}.
}
\]

This solution is always positive and decreases smoothly toward \(0\).

Now apply Euler's method.  Since

\[
F(x,y)=-4y,
\]

the Euler step is

\[
y_{k+1}=y_k+h(-4y_k).
\]

So

\[
\boxed{
y_{k+1}=(1-4h)y_k.
}
\]

If

\[
h=0.6,
\]

then

\[
1-4h=1-2.4=-1.4.
\]

Thus

\[
y_{k+1}=-1.4y_k.
\]

Starting from \(y_0=1\), the Euler values are

\[
1,\quad -1.4,\quad 1.96,\quad -2.744,\quad 3.8416,\ldots
\]

The numerical solution alternates signs and grows in size.  The exact solution does
neither.

\begin{center}
\begin{tabular}{c|c|c|c}
\textbf{step} & \(x_k\) & \textbf{Euler value \(y_k\)} & \textbf{exact value \(e^{-4x_k}\)} \\ \hline
\(0\) & \(0\) & \(1\) & \(1\) \\
\(1\) & \(0.6\) & \(-1.4\) & \(e^{-2.4}\approx0.0907\) \\
\(2\) & \(1.2\) & \(1.96\) & \(e^{-4.8}\approx0.00823\) \\
\(3\) & \(1.8\) & \(-2.744\) & \(e^{-7.2}\approx0.000747\) \\
\(4\) & \(2.4\) & \(3.8416\) & \(e^{-9.6}\approx0.0000677\)
\end{tabular}
\end{center}

The equation says decay.  Euler's method, with this step size, says explosive
oscillation.  The method is lying.

\begin{figure}[h]
\centering
\begin{tikzpicture}
\begin{axis}[
    width=0.84\textwidth,
    height=0.50\textwidth,
    xmin=0, xmax=2.55,
    ymin=-3.2, ymax=4.3,
    axis lines=middle,
    xlabel={\(x\)},
    ylabel={\(y\)},
    xtick={0,0.6,1.2,1.8,2.4},
    ytick={-3,-2,-1,0,1,2,3,4},
    grid=both,
    grid style={line width=.1pt, draw=gray!25},
    legend style={draw=none, at={(0.98,0.97)}, anchor=north east}
]
\addplot[very thick, domain=0:2.5, samples=250] {exp(-4*x)};
\addlegendentry{exact \(e^{-4x}\)}
\addplot[very thick, dashed, mark=*] coordinates {
  (0,1) (0.6,-1.4) (1.2,1.96) (1.8,-2.744) (2.4,3.8416)
};
\addlegendentry{Euler, \(h=0.6\)}
\end{axis}
\end{tikzpicture}
\caption{With a large step, Euler's method can invent behavior opposite to the true solution.}
\label{fig:euler-large-step-false-behavior}
\end{figure}

The problem can be seen from the multiplier

\[
1-4h.
\]

For the exact differential equation, positive solutions decay.  For the Euler values to
decay in size, we need

\[
|1-4h|<1.
\]

This inequality gives

\[
-1<1-4h<1.
\]

Subtract \(1\):

\[
-2<-4h<0.
\]

Divide by \(-4\), reversing the inequalities:

\[
0<h<\frac12.
\]

So for this equation, Euler's method behaves like decay only when

\[
0<h<0.5.
\]

The step \(h=0.6\) is too large.

\begin{quote}
\textbf{Numerical warning.}
A numerical method has its own behavior.  It is not the same object as the
differential equation.

Before trusting a numerical solution, check whether it agrees with the qualitative
behavior of the equation.
\end{quote}

For

\[
y'=-4y,
\]

the qualitative behavior is easy:

\[
y>0 \Longrightarrow y'<0,
\qquad
y<0 \Longrightarrow y'>0.
\]

So solutions are pushed toward \(0\).  A numerical solution that shoots away from
\(0\) is suspect.

This is why numerical methods and direction fields belong together.  The direction field
does not give exact values, but it can catch a numerical approximation that is following
the wrong story.

\subsection{Same differential equation, different initial conditions}
\label{subsec:same-de-different-initial-conditions}

A differential equation gives a law of change.  It does not usually give a single history.

For example,

\[
y'=1-y
\]

has the family of solutions

\[
\boxed{
y(t)=1+Ce^{-t}.
}
\]

Different constants \(C\) give different solution curves.

If

\[
y(0)=0,
\]

then

\[
0=1+C,
\]

so

\[
C=-1,
\]

and

\[
\boxed{
y(t)=1-e^{-t}.
}
\]

This solution starts below \(1\) and rises toward \(1\).

If

\[
y(0)=2,
\]

then

\[
2=1+C,
\]

so

\[
C=1,
\]

and

\[
\boxed{
y(t)=1+e^{-t}.
}
\]

This solution starts above \(1\) and falls toward \(1\).

If

\[
y(0)=1,
\]

then

\[
C=0,
\]

and

\[
\boxed{
y(t)=1.
}
\]

This solution is constant.

The same equation gives all three behaviors.  The initial condition selects which one
happens.

\begin{figure}[h]
\centering
\begin{tikzpicture}
\begin{axis}[
    width=0.84\textwidth,
    height=0.50\textwidth,
    xmin=0, xmax=4.2,
    ymin=-0.2, ymax=2.3,
    axis lines=left,
    xlabel={\(t\)},
    ylabel={\(y(t)\)},
    xtick={0,1,2,3,4},
    ytick={0,1,2},
    grid=both,
    grid style={line width=.1pt, draw=gray!25},
    legend style={draw=none, at={(0.98,0.97)}, anchor=north east}
]
\addplot[very thick, domain=0:4, samples=200] {1-exp(-x)};
\addlegendentry{\(y(0)=0\)}
\addplot[very thick, dashed, domain=0:4, samples=200] {1+exp(-x)};
\addlegendentry{\(y(0)=2\)}
\addplot[very thick, dotted, domain=0:4, samples=2] {1};
\addlegendentry{\(y(0)=1\)}
\addplot[only marks, mark=*] coordinates {(0,0) (0,1) (0,2)};
\end{axis}
\end{tikzpicture}
\caption{The equation \(y'=1-y\) gives a family.  The initial condition selects one solution.}
\label{fig:same-de-different-initial-conditions}
\end{figure}

For a first-order equation, one initial condition usually chooses one solution curve.
For a second-order equation, one initial condition is usually not enough.

Consider the oscillator equation

\[
y''+y=0.
\]

The general solution is

\[
y(t)=C_1\cos t+C_2\sin t.
\]

If we know only

\[
y(0)=1,
\]

then

\[
1=C_1.
\]

But \(C_2\) is still free.  Both functions

\[
y_1(t)=\cos t
\]

and

\[
y_2(t)=\cos t+2\sin t
\]

solve

\[
y''+y=0
\]

and both satisfy

\[
y(0)=1.
\]

They differ because their initial velocities differ:

\[
y_1'(0)=0,
\qquad
y_2'(0)=2.
\]

A second-order motion needs position and velocity.

\begin{quote}
\textbf{Initial data rule.}
A differential equation gives the rule of motion.  Initial conditions give the starting
state.

A first-order equation usually needs one starting value.  A second-order equation
usually needs two: position and velocity.
\end{quote}

This is why a spring cannot be described by its position alone.  If the mass is at
equilibrium, it may stay there, or it may be passing through with high speed.  The
position is the same.  The future is not.

\subsection{The three warnings together}
\label{subsec:three-warnings-together}

The warnings can be remembered as three questions.

\begin{center}
\begin{tabular}{p{0.30\textwidth}|p{0.58\textwidth}}
\textbf{Before trusting the result, ask} & \textbf{What can go wrong?} \\ \hline
Did we divide by an expression that might be zero? &
We may have lost equilibrium solutions. \\[4pt]
Is a numerical step size too large? &
Euler's method may create false oscillation or false growth. \\[4pt]
Have we supplied enough initial information? &
The differential equation may describe a family, not one solution.
\end{tabular}
\end{center}

Differential equations are powerful because they describe change locally.  That same
local nature makes them sensitive to algebraic assumptions, numerical choices, and
initial data.

The chapter ends by gathering the main methods and putting them back to work in
applications.

\section{Chapter review and applications}
\label{sec:de-chapter-review-applications}

This chapter began with a sentence:

\[
\text{A differential equation is a law written in the language of change.}
\]

We have now seen several versions of that sentence.

A cooling cup obeys a law involving temperature difference.  A population obeys a law
involving current population and available capacity.  A tank obeys a law involving inflow
and outflow.  A spring obeys a law involving displacement, velocity, and acceleration.
Two populations can obey linked laws.

The methods differ, but the pattern is the same:

\[
\boxed{
\text{model the rate, then understand the function.}
}
\]

\subsection{Skills: solve and interpret first-order equations}
\label{subsec:skills-solve-interpret-first-order-equations}

The first task in a differential equation is not to solve.  It is to recognize what kind of
equation is in front of us and what kind of answer is reasonable.

\begin{center}
\begin{tabular}{p{0.24\textwidth}|p{0.29\textwidth}|p{0.34\textwidth}}
\textbf{form} & \textbf{main method} & \textbf{main question it answers} \\ \hline
\(y'=F(t,y)\) & direction field or Euler's method & What do solution curves seem to do? \\[4pt]
\(y'=f(y)\) & phase line & Where are equilibria, and are they stable? \\[4pt]
\(y'=g(t)h(y)\) & separation & Can the variables be separated and integrated? \\[4pt]
\(y'+p(t)y=q(t)\) & integrating factor & Can the product rule be used backward? \\[4pt]
\(P'=rP(1-P/K)\) & logistic formula and phase line & How does growth slow near carrying capacity? \\[4pt]
\(Q'=-k(Q-Q_{\text{eq}})\) & exponential approach to equilibrium & How does a quantity approach a stable level?
\end{tabular}
\end{center}

The following examples review the central skills.

\paragraph{Example 1: separable growth.}
Solve

\[
y'=3ty,
\qquad
y(0)=2.
\]

Separate:

\[
\frac{dy}{y}=3t\,dt.
\]

Integrate:

\[
\ln|y|=\frac32t^2+C.
\]

Since \(y(0)=2\),

\[
C=\ln2.
\]

Therefore

\[
\boxed{
y(t)=2e^{(3/2)t^2}.
}
\]

The exponent is the accumulated relative growth rate:

\[
\int_0^t 3s\,ds=\frac32t^2.
\]

\paragraph{Example 2: cooling toward equilibrium.}
Solve

\[
T'=-0.1(T-22),
\qquad
T(0)=82.
\]

This has the equilibrium-difference form

\[
Q'=-k(Q-Q_{\text{eq}}).
\]

So

\[
T(t)=22+(82-22)e^{-0.1t}.
\]

Thus

\[
\boxed{
T(t)=22+60e^{-0.1t}.
}
\]

The temperature approaches \(22^\circ\text{C}\), not \(0^\circ\text{C}\), because \(22\) is
the surrounding temperature.

\paragraph{Example 3: a first-order linear equation.}
Solve

\[
y'+2y=4e^{-t},
\qquad
y(0)=3.
\]

Here

\[
p(t)=2,
\qquad
q(t)=4e^{-t}.
\]

An integrating factor is

\[
\mu(t)=e^{\int 2\,dt}=e^{2t}.
\]

Multiply the equation by \(e^{2t}\):

\[
e^{2t}y'+2e^{2t}y=4e^t.
\]

The left side is

\[
(e^{2t}y)'.
\]

Therefore

\[
(e^{2t}y)'=4e^t.
\]

Integrate:

\[
e^{2t}y=4e^t+C.
\]

So

\[
y=4e^{-t}+Ce^{-2t}.
\]

Use \(y(0)=3\):

\[
3=4+C.
\]

Thus

\[
C=-1.
\]

Therefore

\[
\boxed{
y(t)=4e^{-t}-e^{-2t}.
}
\]

\paragraph{Example 4: logistic interpretation.}
Suppose

\[
P'=0.5P\left(1-\frac{P}{800}\right).
\]

The equilibria are

\[
P=0
\qquad\text{and}\qquad
P=800.
\]

If

\[
0<P<800,
\]

then

\[
P'>0.
\]

If

\[
P>800,
\]

then

\[
P'<0.
\]

So positive solutions move toward

\[
P=800.
\]

The number \(800\) is the carrying capacity.  The equilibrium \(P=800\) is stable, while
\(P=0\) is unstable from the right.

\begin{quote}
\textbf{Review habit.}
After solving a differential equation, ask:

\begin{enumerate}
\item Does the solution satisfy the differential equation?

\item Does it satisfy the initial condition?

\item Does its long-term behavior match the model?

\item Do the units make sense?

\item Were any solutions lost by division?
\end{enumerate}
\end{quote}

\subsection{Concept check}
\label{subsec:de-concept-check}

These questions are meant to test meanings before procedures.

\begin{enumerate}
\item What is the difference between a differential equation and a solution of a
differential equation?

\item Why does the equation

\[
y'=1-y
\]

suggest that solutions should move toward \(y=1\)?

\item What does a direction field draw at each point \((t,y)\)?

\item What is an equilibrium solution of an autonomous equation

\[
y'=f(y)?
\]

\item Why can a phase line summarize an autonomous equation but not a general
equation

\[
y'=F(t,y)?
\]

\item Why does Euler's method usually improve when the step size is smaller?

\item Why is smaller step size not the same thing as a proof of accuracy?

\item When separating variables, why must we check constant solutions separately?

\item What is the difference between a stable equilibrium and an unstable equilibrium?

\item Why does a second-order equation usually require two initial conditions?
\end{enumerate}

\subsection{Skills practice}
\label{subsec:de-skills-practice}

Solve or analyze each problem.

\begin{enumerate}
\item Verify that

\[
y(t)=2e^{-3t}
\]

solves

\[
y'=-3y,
\qquad
y(0)=2.
\]

\item Solve

\[
y'=5y,
\qquad
y(0)=4.
\]

\item Solve

\[
y'=t^2y,
\qquad
y(0)=1.
\]

\item Solve

\[
y'=\frac{x}{1+y^2},
\qquad
y(0)=0,
\]

leaving the answer implicitly if needed.

\item Solve

\[
T'=-0.04(T-18),
\qquad
T(0)=78.
\]

Then find the limiting temperature as \(t\to\infty\).

\item A tank contains \(200\) liters of water.  Brine with concentration \(0.3\) grams per
liter enters at \(4\) liters per minute, and the well-mixed solution leaves at the same
rate.  Let \(A(t)\) be the grams of salt in the tank.  Set up and solve the initial value
problem if the tank initially contains pure water.

\item Solve the linear equation

\[
y'+3y=6,
\qquad
y(0)=10.
\]

\item Solve

\[
y'-2y=e^t,
\qquad
y(0)=0.
\]

\item For

\[
y'=y(y-2),
\]

find the equilibria and draw a phase line.  Classify each equilibrium as stable or
unstable.

\item For

\[
P'=0.2P\left(1-\frac{P}{500}\right),
\qquad
P(0)=50,
\]

write the logistic solution.

\item Use Euler's method with \(h=0.25\) to approximate \(y(1)\) for

\[
y'=1-y,
\qquad
y(0)=0.
\]

\item Use improved Euler's method with \(h=0.5\) to approximate \(y(1)\) for

\[
y'=1-y,
\qquad
y(0)=0.
\]

\item Check that

\[
y(t)=3\cos(2t)-\sin(2t)
\]

solves

\[
y''+4y=0.
\]

Find \(y(0)\) and \(y'(0)\).

\item Rewrite

\[
y''+9y=0
\]

as a first-order system using \(v=y'\).
\end{enumerate}

\subsection{Interpretation and modeling practice}
\label{subsec:de-interpretation-modeling-practice}

\begin{enumerate}
\item The equation

\[
M'=-0.02M
\]

models the mass of a radioactive substance.  Explain in words what the equation says.
What are the units of \(0.02\) if time is measured in years?

\item The equation

\[
T'=-k(T-T_s)
\]

has \(k>0\).  Explain why the sign of \(T'\) is correct both when \(T>T_s\) and when
\(T<T_s\).

\item A population satisfies

\[
P'=0.1P\left(1-\frac{P}{1000}\right).
\]

Explain why the model predicts nearly exponential growth when \(P\) is small compared
with \(1000\).

\item A population satisfies the threshold model

\[
P'=0.01P(100-P)(P-20).
\]

Find the equilibria and explain what happens if

\[
P(0)=15
\]

and if

\[
P(0)=25.
\]

\item A fish population satisfies

\[
P'=0.4P\left(1-\frac{P}{1000}\right)-h.
\]

Find the largest sustainable constant harvest rate in this model.

\item The solution of a cooling model is

\[
T(t)=20+70e^{-0.05t}.
\]

What are the surrounding temperature, the initial temperature, and the cooling
constant?

\item A bank account balance satisfies

\[
B'=0.03B+1000.
\]

Explain the meaning of each term.  What are the units of \(1000\)?

\item A numerical solution of

\[
y'=-5y,
\qquad
y(0)=1
\]

produces values that alternate sign and grow in size.  Explain why this contradicts
the direction field.
\end{enumerate}

\subsection{Project: model a cooling cup of coffee}
\label{subsec:project-cooling-cup-coffee}

Newton's law of cooling says

\[
T'=-k(T-T_s),
\]

where \(T_s\) is the surrounding temperature and \(k>0\).

The solution is

\[
\boxed{
T(t)=T_s+(T_0-T_s)e^{-kt}.
}
\]

This project asks you to test that model with data.

\paragraph{Data collection.}
Use a hot drink, soup, or warm water in a cup.

\begin{enumerate}
\item Measure the room temperature \(T_s\).

\item Measure the initial temperature \(T_0\).

\item Record the temperature every \(2\) or \(3\) minutes for at least \(20\) minutes.

\item Make a table of \(t\), \(T(t)\), and \(T(t)-T_s\).
\end{enumerate}

\paragraph{Linearizing the model.}
If the cooling model is correct, then

\[
T(t)-T_s=(T_0-T_s)e^{-kt}.
\]

Assuming \(T(t)>T_s\), take logarithms:

\[
\ln(T(t)-T_s)=\ln(T_0-T_s)-kt.
\]

So a plot of

\[
\ln(T(t)-T_s)
\]

against \(t\) should be approximately a straight line with slope

\[
-k.
\]

\paragraph{Project tasks.}

\begin{enumerate}
\item Plot \(T(t)\) against \(t\).  Does the curve look like exponential approach to room
temperature?

\item Plot \(\ln(T(t)-T_s)\) against \(t\).  Does the graph look approximately linear?

\item Estimate \(k\) from the slope of the logarithmic plot.

\item Use your estimated \(k\) to predict when the drink will reach a chosen temperature,
such as \(45^\circ\text{C}\).

\item Compare your prediction with the observed data.

\item Describe at least two possible reasons why real data might not perfectly follow the
model.
\end{enumerate}

\begin{quote}
\textbf{Model reflection.}
A differential equation is a simplified law.  Real cups lose heat through the cup, the
surface, evaporation, stirring, and changing room conditions.  The model is useful
because it captures the main pattern, not because it records every physical detail.
\end{quote}

\subsection{Project: compare Euler approximations with exact solutions}
\label{subsec:project-euler-compare-exact}

Euler's method turns a differential equation into a sequence:

\[
y_{k+1}=y_k+hF(x_k,y_k).
\]

This project compares Euler approximations with exact solutions.

\paragraph{Part A: a stable approach.}
Use

\[
y'=1-y,
\qquad
y(0)=0.
\]

The exact solution is

\[
y=1-e^{-x}.
\]

For each step size

\[
h=0.5,\qquad h=0.25,\qquad h=0.1,
\]

use Euler's method to approximate \(y(2)\).

Make a table:

\begin{center}
\begin{tabular}{c|c|c|c}
\textbf{step size} & \textbf{Euler approximation} & \textbf{exact value} & \textbf{absolute error} \\ \hline
\(0.5\) & & & \\
\(0.25\) & & & \\
\(0.1\) & & &
\end{tabular}
\end{center}

Then answer:

\begin{enumerate}
\item Does the error decrease as \(h\) decreases?

\item Does Euler's method overestimate or underestimate the exact solution?

\item How does the concavity of \(y=1-e^{-x}\) help explain the sign of the error?
\end{enumerate}

\paragraph{Part B: a step-size warning.}
Use

\[
y'=-4y,
\qquad
y(0)=1.
\]

The exact solution is

\[
y=e^{-4x}.
\]

Use Euler's method up to \(x=2.4\) with

\[
h=0.2
\qquad\text{and}\qquad
h=0.6.
\]

Then answer:

\begin{enumerate}
\item Which step size gives values that behave like the exact solution?

\item Which step size gives values that alternate sign?

\item For this equation, Euler's multiplier is

\[
1-4h.
\]

Explain why the condition

\[
|1-4h|<1
\]

is related to numerical decay.
\end{enumerate}

\paragraph{Part C: optional improved Euler.}
Repeat Part A with improved Euler's method:

\[
k_1=F(x_k,y_k),
\]

\[
\widetilde y_{k+1}=y_k+hk_1,
\]

\[
k_2=F(x_k+h,\widetilde y_{k+1}),
\]

\[
y_{k+1}=y_k+\frac h2(k_1+k_2).
\]

Compare its error with ordinary Euler's method for the same step sizes.

\begin{quote}
\textbf{Project conclusion.}
A numerical method should be judged by both computation and behavior.  A table of
numbers is not enough.  The numbers should agree with the direction field, the phase
line, and any exact solution available.
\end{quote}

\subsection{A final map of the chapter}
\label{subsec:final-map-differential-equations}

The main ideas of the chapter can be summarized in one picture of dependencies.

\[
\text{law of change}
\quad\Longrightarrow\quad
\text{differential equation}
\]

and then we choose a route:

\[
\begin{aligned}
\text{qualitative route} &:\quad \text{direction field, phase line, equilibria},\\
\text{numerical route} &:\quad \text{Euler's method and step-size checks},\\
\text{exact route} &:\quad \text{separation or integrating factors},\\
\text{modeling route} &:\quad \text{interpret constants, units, and long-term behavior}.
\end{aligned}
\]

No one route is the whole subject.  Often the best solution uses several at once.

For example, in a logistic model, the phase line tells us that a positive solution should
approach \(K\).  The exact formula tells us when it reaches \(K/2\).  Euler's method
can approximate values when parameters are changing or when the equation becomes too
complicated to solve exactly.  The model interpretation tells us which values of \(P\)
make physical sense.

\begin{quote}
\textbf{What differential equations add to calculus.}
Earlier, calculus measured change and accumulated change.  Differential equations use
those ideas to describe laws.

They turn

\[
\text{the derivative}
\]

from something we compute into something we are given.
\end{quote}

This chapter ends one story and opens another.

Differential equations often describe motion by using time as the input:

\[
t\longmapsto y(t).
\]

But many curves are not naturally graphs of the form

\[
y=f(x).
\]

A point on a wheel, a planet moving in an orbit, or a particle tracing a loop may pass
above the same \(x\)-value more than once.  The curve may still be easy to describe if
we let time trace both coordinates:

\[
x=x(t),
\qquad
y=y(t).
\]

That is the next idea.  Instead of forcing every curve to be a graph, we let a moving
point draw it.
